%%% Hlavní soubor. Zde se definují základní parametry a odkazuje se na ostatní části. %%%

%% Verze pro jednostranný tisk:
% Okraje: levý 40mm, pravý 25mm, horní a dolní 25mm
% (ale pozor, LaTeX si sám přidává 1in)
\documentclass[12pt,a4paper]{report}
\setlength\textwidth{145mm}
\setlength\textheight{247mm}
\setlength\oddsidemargin{15mm}
\setlength\evensidemargin{15mm}
\setlength\topmargin{0mm}
\setlength\headsep{0mm}
\setlength\headheight{0mm}
% \openright zařídí, aby následující text začínal na pravé straně knihy
\let\openright=\clearpage

\renewcommand\baselinestretch{1.3} % riadkovanie jeden a pol

%% Pokud tiskneme oboustranně:
% \documentclass[12pt,a4paper,twoside,openright]{report}
% \setlength\textwidth{145mm}
% \setlength\textheight{247mm}
% \setlength\oddsidemargin{15mm}
% \setlength\evensidemargin{0mm}
% \setlength\topmargin{0mm}
% \setlength\headsep{0mm}
% \setlength\headheight{0mm}
% \let\openright=\cleardoublepage

%% Pokud pouľíváte csLaTeX (doporučeno):
%\usepackage{czech}
%% Pokud nikoliv:
%\usepackage[czech]{babel}
%\usepackage[T1]{fontenc}

%% Pouľité kódování znaků: obvykle latin2, cp1250 nebo utf8:
\usepackage[utf8]{inputenc}
\usepackage[T1]{fontenc}

%% Ostatní balíčky
\usepackage{graphicx}
\usepackage{amsthm}
\usepackage{amssymb}
\usepackage{alltt}
\usepackage{todonotes}
\usepackage{algorithmic}
\usepackage[chapter]{algorithm}
\usepackage{enumerate}
\usepackage{multirow}
\usepackage{threeparttable}
\usepackage{listings}
\usepackage{algorithmic}
\usepackage[chapter]{algorithm}

\usepackage{color}
\definecolor{dkgreen}{rgb}{0,0.6,0}
\definecolor{gray}{rgb}{0.5,0.5,0.5}
\definecolor{mauve}{rgb}{0.58,0,0.82}
\newcommand{\lstfont}{\small\ttfamily}
\lstset{ %
  language=,                % the language of the code
  basicstyle=\lstfont,           % the size of the fonts that are used for the code
  %numbers=left,                   % where to put the line-numbers
  %numberstyle=\tiny\color{gray},  % the style that is used for the line-numbers
  %stepnumber=2,                   % the step between two line-numbers. If it's 1, each line 
                                  % will be numbered
  numbersep=5pt,                  % how far the line-numbers are from the code
  backgroundcolor=\color{white},      % choose the background color. You must add \usepackage{color}
  showspaces=false,               % show spaces adding particular underscores
  showstringspaces=false,         % underline spaces within strings
  showtabs=false,                 % show tabs within strings adding particular underscores
  frame=single,                   % adds a frame around the code
  rulecolor=\color{black},        % if not set, the frame-color may be changed on line-breaks within not-black text (e.g. commens (green here))
  tabsize=2,                      % sets default tabsize to 2 spaces
  captionpos=b,                   % sets the caption-position to bottom
  breaklines=true,                % sets automatic line breaking
  breakatwhitespace=false,        % sets if automatic breaks should only happen at whitespace
  title=\lstname,                   % show the filename of files included with \lstinputlisting;
                                  % also try caption instead of title
  keywordstyle=\color{blue},          % keyword style
  commentstyle=\color{dkgreen},       % comment style
  stringstyle=\color{mauve},         % string literal style
  escapeinside={\%*}{*)},            % if you want to add a comment within your code
  morekeywords={*,...},              % if you want to add more keywords to the set
}

\renewcommand{\algorithmicrequire}{\textbf{Input:}}
\renewcommand{\algorithmicensure}{\textbf{Output:}}

\renewcommand{\algorithmicforall}{\textbf{for each}}

\renewcommand{\algorithmiccomment}[1]{// #1}

% pekne pokope definujeme potrebne udaje
\def\mftitle{Inference of an XML Schema with the Knowledge of XML Operations}
\def\mfthesistype{Master Thesis}
\def\mfauthor{Mário Mikula}
\def\mfadvisor{RNDr. Irena Mlýnková, Ph.D.}
\def\mfplacedate{Praha, 2011}

%% Balíček hyperref, kterým jdou vyrábět klikací odkazy v PDF,
%% ale hlavně ho pouľíváme k uloľení metadat do PDF (včetně obsahu).
%% POZOR, nezapomeňte vyplnit jméno práce a autora.
\usepackage[ps2pdf,unicode]{hyperref}   % Musí být za vąemi ostatními balíčky
\hypersetup{pdftitle=\mftitle}
\hypersetup{pdfauthor=\mfauthor}

%%% Drobné úpravy stylu

% Tato makra přesvědčují mírně oąklivým trikem LaTeX, aby hlavičky kapitol
% sázel příčetněji a nevynechával nad nimi spoustu místa. Směle ignorujte.
%\makeatletter
%\def\@makechapterhead#1{
%  {\parindent \z@ \raggedright \normalfont
%   \Huge\bfseries \thechapter. #1
%   \par\nobreak
%   \vskip 20\p@
%}}
%\def\@makeschapterhead#1{
%  {\parindent \z@ \raggedright \normalfont
%   \Huge\bfseries #1
%   \par\nobreak
%   \vskip 20\p@
%}}
%\makeatother

% Toto makro definuje kapitolu, která není očíslovaná, ale je uvedena v obsahu.
\def\chapwithtoc#1{
\chapter*{#1}
\addcontentsline{toc}{chapter}{#1}
}


% Theorem environment
\newenvironment{definition}[1][Definition]{\begin{trivlist}
\item[\hskip \labelsep {\bfseries #1}]}{\end{trivlist}}


\begin{document}

% Trochu volnějąí nastavení dělení slov, neľ je default.
\lefthyphenmin=2
\righthyphenmin=2

%%% Titulní strana práce

\pagestyle{empty}
\begin{center}

\large

Charles University in Prague

\medskip

Faculty of Mathematics and Physics

\vfill

{\bf\Large \mfthesistype}

\vfill

\centerline{\mbox{\includegraphics[width=60mm]{logo.eps}}}

\vfill
\vspace{5mm}

{\LARGE \mfauthor}

\vspace{15mm}

% Název práce přesně podle zadání
{\LARGE\bfseries \mftitle}

\vfill

% Název katedry nebo ústavu, kde byla práce oficiálně zadána
% (dle Organizační struktury MFF UK)
Department of Software Engineering

\vfill

\begin{tabular}{rl}

Supervisor of the master thesis: & \mfadvisor \\
\noalign{\vspace{2mm}}
Study programme: & Informatika \\
\noalign{\vspace{2mm}}
Specialization: & ISS \\
\end{tabular}

\vfill

% Zde doplňte rok
\mfplacedate

\end{center}

\newpage

%%% Následuje vevázaný list -- kopie podepsaného "Zadání diplomové práce".
%%% Toto zadání NENÍ součástí elektronické verze práce, nescanovat.

%%% Na tomto místě mohou být napsána případná poděkování (vedoucímu práce,
%%% konzultantovi, tomu, kdo zapůjčil software, literaturu apod.)

\openright

\noindent
I would like to thank my supervisor RNDr. Irena Mlýnková, Ph.D., for her helpful advices, corrections and suggestions.

\newpage

%%% Strana s čestným prohláąením k diplomové práci

\vglue 0pt plus 1fill

\noindent
I declare that I carried out this master thesis independently, and only with the cited
sources, literature and other professional sources.

\medskip\noindent
I understand that my work relates to the rights and obligations under the Act No.
121/2000 Coll., the Copyright Act, as amended, in particular the fact that the Charles
University in Prague has the right to conclude a license agreement on the use of this
work as a school work pursuant to Section 60 paragraph 1 of the Copyright Act.

\vspace{10mm}

\hbox{\hbox to 0.5\hsize{%
In ........ date ............
\hss}\hbox to 0.5\hsize{%
Signature
\hss}}

\vspace{20mm}
\newpage

%%% Povinná informační strana diplomové práce

\vbox to 0.5\vsize{
\setlength\parindent{0mm}
\setlength\parskip{5mm}

Název práce:
Odvozování XML schémy se znalostí XML operací
% přesně dle zadání

Autor:
\mfauthor

Katedra:  % Případně Ústav:
Katedra softwarového inženýrství
% dle Organizační struktury MFF UK

Vedoucí diplomové práce:
\mfadvisor
% dle Organizační struktury MFF UK, případně plný název pracoviątě mimo MFF UK

Abstrakt: V poslednej dobe bolo vyvinutých množstvo metód zaoberajúcich sa automatickým odvodzovaním XML schémy. Väčšina z nich ale používa XML dokumenty ako jediný vstup. My sa v tejto práci zameriavame na rozšírenie odvo\-dzo\-vania zahrnutím XML operácií, konkrétne XQuery dotazov. Diskutujeme ako je možné využiť XQuery dotazy na vylepšenie procesu odvodzovania a navrhujeme algoritmus založený na vybraných vylepšeniach, ktorý rozširuje existujúcu metódu hľadania kľúčov a ktorý môže byť začlenený do tzv. \emph{metód odvodzujúcich gramatiku}. Implementovaním algoritmu sme vytvorili prvé kompletné riešenie odvodzovania XML schémy, ktoré využíva XML dokumenty spolu s XML operáciami.
% abstrakt v rozsahu 80-200 slov; nejedná se vąak o opis zadání diplomové práce

Klíčová slova: XML, XML schéma, odvozování, XQuery, dotazy, XML operace
% 3 aľ 5 klíčových slov

\vss}\nobreak\vbox to 0.49\vsize{
\setlength\parindent{0mm}
\setlength\parskip{5mm}

Title:
\mftitle
% přesný překlad názvu práce v angličtině

Author:
\mfauthor

Department:
Department of Software Engineering
% dle Organizační struktury MFF UK v angličtině

Supervisor:
\mfadvisor
% dle Organizační struktury MFF UK, případně plný název pracoviątě
% mimo MFF UK v angličtině

Abstract: Recently, plenty of methods dealing with automatic inference of XML schema have been developed, however, most of them utilize XML documents as their only input. In this thesis we focus on extending inference by incorporating XML operations, in particular XQuery queries. We discuss how can XQuery queries help in improving the inference process and we propose an algorithm based on chosen improvements, extending an existing method of a key discovery, that can be integrated to so-called \emph{grammar-inferring} methods. By implementing it, we created the first solution of XML schema inference using XML documents along with XML operations.
% abstrakt v rozsahu 80-200 slov v angličtině; nejedná se vąak o překlad
% zadání diplomové práce 

Keywords: XML, XML schema, inference, XQuery, queries, XML operations
% 3 aľ 5 klíčových slov v angličtině

\vss}

\newpage

%%% Strana s automaticky generovaným obsahem diplomové práce. U matematických
%%% prací je přípustné, aby seznam tabulek a zkratek, existují-li, byl umístěn
%%% na začátku práce, místo na jejím konci.

\openright
\pagestyle{plain}
\renewcommand\thepage{}
\tableofcontents

\newtheorem{theorem}{Theorem}
\theoremstyle{definition}
%\newtheorem*{define}{Definition}	% Definice nečíslujeme, proto "*"
\newtheorem{define}{Definition}[chapter]
\newtheorem{example}{Example}

\newpage
\renewcommand\thepage{\arabic{page}}
\setcounter{page}{1}

%%% Jednotlivé kapitoly práce jsou pro přehlednost uloľeny v samostatných souborech
\chapter{Introduction}
The Extensible Markup Language (XML) \cite{Bray:08:EML} has become a popular standard for data representation and exchange. Its prevalence is based on its simplicity and, at the same time, its flexibility and expressive power.

With XML used to represent data themselves, it is often needed or convenient to somehow specify the structure of the represented data, their format, internal relations or restrictions, etc. In order to this need, several so-called XML schema languages (XML schemata) have been created. The most used of them are Document Type Definition (DTD) \cite{Bray:08:EML} and XML Schema (or XSD as XML Schema Definition) \cite{Walmsley:04:XSP, Thompson:04:XSP, Malhotra:04:XSP} proposed by the World Wide Web Consorcium (W3C).

Despite a conveyed encouragement to use an XML schema along with an XML data representation, in practise it is done sparsely. Commonly, XML documents are not assigned with their respective XML schema at all or the schema is outdated due to modifications done to a structure of data without updating the schema \cite{Mlynkova:2008:AAX:1494650.1495496}.

Recently in reaction to this situation, a significant number of approaches dealing with automatic construction of an XML schema has been proposed. The aim is to exploit a provided set of XML documents and infer an XML schema, so that the XML documents are valid against it. In addition, the inferred schema should be reasonable in views of human-readability, preciseness and conciseness. Most of published approaches to XML schema inference are of this type - the input is a set of XML documents. They are based on various ideas and can be classified by several aspects as discussed in Chapter \ref{chapter_analysis_of_recent_approaches}.

Besides the mentioned type, some approaches that utilize other or additional sources have been developed, for example \cite{Mlynkova:2009:IXS:1862681.1862693} and \cite{thesis_klempa}. If there are available sources like an out-dated XML schema, operations upon the XML data such as a set of XQuery \cite{w3c_xquery} queries, or other, these sources can be exploited to refine the process of inference. The refinement can be achieved in various aspects such as increasing the speed of the process, getting a more precise, more concise or more readable result or inference of statements about the data which cannot be (easily) extracted from the data themselves.

Recently, the main effort has been focused on a research of the approaches that utilize XML documents, and thus, there are only few approaches of the latter type (as also discussed in Chapter \ref{chapter_analysis_of_recent_approaches}), leaving a wide space for a possible future research and improvements.

\section{Aim of this Work}
Our aim is to perform a research on the problem of inference of an XML schema for the given set of XML data in a situation when we are provided also with a set of related operations (XML queries, XSLT scripts \cite{Clark:99:XTV} etc.).

Firstly, we analyze existing inference solutions in general and discuss their advantages and disadvantages. Then, we identify and discuss information that can be extracted from a given set of XML operations and how they can be exploited to achieve more precise and realistic XML schema.

The work results in a proposal of own approach involving the improvements, its implementation and suitable experiments that show its advantages.

\section{Structure of the Thesis}
The thesis is structured as follows.

In Chapter 2, we introduce XML technologies used in this thesis and we briefly explain their basic principles.

Chapter 3 contains summaries of several existing approaches of XML schema inference. One of them proposes an utilisation of XQuery queries, which interests us the most.

We made an analysis of XQuery technology from a view of XML schema inference, where we suggest several ideas on how XQuery queries can be exploited in the inference process. This is a content of Chapter 4.

In Chapter 5, we discuss some essential questions that we needed to decide before a development of an algorithm.

The algorithm itself is proposed in Chapter 6. And Chapter 7 describes how to combine results of the algorithm with existing grammar inferring methods of XML schema inference.

An implementation of the proposed algorithm using the jInfer framework is described in Chapter 8. Chapter 9 then discusses performed experiments and their outcomes. And, finally, Chapter 10 concludes.

\chapter{Used Technologies and Definitions}

\section{XML Schema}
An XML schema refers to a description of an XML document in terms of its structure and various constraints. Commonly, the XML schema describes element and attribute names, their parent-child relations, their order and a type of their content. Other constraints often expressed in the XML schema are restrictions on numbers of occurrences of elements, specification of (non-)obligatory attributes, uniqueness and specification of keys.

Various languages have been proposed to express XML schemata. The most known are Document Type Definition (DTD) \cite{Bray:08:EML} and XML Schema Definition (XML Schema, XSD) \cite{Walmsley:04:XSP, Thompson:04:XSP, Malhotra:04:XSP} which are briefly described in the following sections. Another examples of the XML schema languages are RELAX NG \cite{relaxng} and  Schematron \cite{schematron}.

Validity of an XML document against its XML schema expresses whether the document is well-formed \cite{Bray:08:EML} and, at the same time, whether it conforms to the XML schema.

\subsection{An XML Example}
To demonstrate the described technologies, we introduce a part of a simple XML document (see Figure \ref{FIG_a_simple_xml_example}). It represents books in a bookstore. Each book has a mandatory id, a title, a list of authors and an optional ISBN.

\begin{figure}
\caption{A simple XML example.}
\label{FIG_a_simple_xml_example}
\begin{verbatim}
<bookstore>
  <book id="b1">
    <title>
      Compilers: Principles, Techniques, and Tools
      (2nd Edition)
    </title>
    <authors>
      <author>Alfred V. Aho</author>
      <author>Monica S. Lam</author>
      <author>Ravi Sethi</author>
      <author>Jeffrey D. Ullman</author>
    </authors>
  </book>
  <book id="b2">
    <title>XQuery</title>
    <authors>
      <author>Priscilla Walmsley</author>
    </authors>
    <isbn>0596006349</isbn>
  </book>
</bookstore>
\end{verbatim}
\end{figure}

\subsection{DTD}
Document Type Definition (DTD) expresses a structure of XML documents by declarations of elements. An element has its name and a content declared using syntax \texttt{<!ELEMENT $name$ $content$>}.

The content of an element can be denoted by \texttt{EMPTY} for an empty element, \texttt{ANY} for any content, \texttt{(\#PCDATA)} allowing only textual content (without any other subelements), or specified using regular expressions. Names of subelements are combined using operators (\texttt{|}, \texttt{+}, \texttt{*}, \texttt{?} and \texttt{,}(comma)). To express the mixed content \texttt{\#PCDATA} can be used in an alternation list with the subelement names and this alternation has to be enclosed in \texttt{*} operator.

Attributes of an element are specified in an attribute list \texttt{<!ATTLIST\linebreak $element\_name$ $attribute\_name$ $type$ $default\_value$>}. Each attribute has its name, its type and its default value or definition of obligation of occurrence. The type can be an enumeration of values \texttt{($value_1$ | $value_2$ | $\dots$ | $value_n$)} or one of the following values.

\begin{itemize}
\item \texttt{CDATA} Character data - any string.
\item \texttt{ID} An unique identifier.
\item \texttt{IDREF} An ID reference - a value of the ID.
\item \texttt{IDREFS} A space-separated list of ID references. 
\item \texttt{NMTOKEN}	A valid XML name.
\item \texttt{NMTOKENS} A space-separated list of valid XML names.
\item \texttt{ENTITY} An entity.
\item \texttt{ENTITIES} A space-separated list of entities.
\item \texttt{NOTATION} A name of a notation.
\end{itemize}

The dafault value is either a literal value or one of the following specifiers.

\begin{itemize}
\item \texttt{\#REQUIRED} The attribute is mandatory.
\item \texttt{\#IMPLIED} The attribute is optional.
\item \texttt{\#FIXED $value$} The attribute value is constant $value$.
\end{itemize}

The DTD also provides other constructs such as declaration of entities and notations not mentioned in this work.

An example of DTD describing the book element from the XML example in Figure \ref{FIG_a_simple_xml_example} follows.

\begin{verbatim}
<!ELEMENT book (title, authors, isbn?)>
<!ATTRLIST book id ID #REQUIRED>

<!ELEMENT title #PCDATA>

<!ELEMENT authors (author+)>

<!ELEMENT author #PCDATA>

<!ELEMENT isbn #PCDATA>
\end{verbatim}

\subsection{XSD}
Since the XSD language, containing many constructs and features, is quite comprehensive, we will describe just its basic principles. An important fact is that each XSD instance is a valid XML document. Its root element is always \texttt{<schema>} and the XSD instance begins with the following two lines.

\begin{verbatim}
<?xml version="1.0" encoding="UTF-8"?>
<xs:schema xmlns:xs="http://www.w3.org/2001/XMLSchema">
\end{verbatim}

Definition of an element in the XSD is
\begin{alltt}
<xs:element name="\(name\)" type="\(type\)"/>
\end{alltt}
where $type$ is either one of the built-in types or a user-defined type. Definition of an attribute is similar.
\begin{alltt}
<xs:attribute name="\(name\)" type="\(type\)"/>
\end{alltt}
Attributes are optional by default. If an attribute is mandatory, it is expressed by adding \texttt{use} attribute to its definition.
\begin{alltt}
<xs:attribute name="\(name\)" type="\(type\)" use="required"/>
\end{alltt}

The user-defined type in the XSD language is either a simple type, if it does not contain other elements and attributes, or a complex type otherwise. Attributes are only allowed to be of the simple types. Definition of a simple type is
\begin{alltt}
<xs:simpleType name="\(name\)">
   \(type details\)
</xs:simpleType>
\end{alltt}
The user-defined simple types often serve to restrict the built-in types in various ways such as limiting lengths of strings, allowing only certain values and thus creating an enumeration type, and other.

A complex type is defined in the same way.
\begin{alltt}
<xs:complexType name="\(name\)">
  \(type details\)
</xs:complexType>
\end{alltt}
The complex types can contain many constructs. Subelements are declared using \texttt{<xs:sequence>}, \texttt{<xs:choice>} and \texttt{<xs:all>} schema elements. If the order of the subelements is significant, we use \texttt{<xs:sequence>}, where the subelements in an XML instance must occur in the same order as they are defined in the sequence. If the order is not significant, we use \texttt{<xs:all>}. Construct \texttt{<xs:choice>} is equivalent to the alternation of several elements in the DTD.

Moreover, these three construct can be nested and combined, assuming the combination is not ambiguous.

Many schema elements (including \texttt{<xs:element>}, \texttt{<xs:sequence>}, \texttt{<xs:choice>}, \texttt{<xs:all>}) can be assigned with an occurrence interval. The occurrence interval is expressed using \texttt{minOccurs} and \texttt{maxOccurs} attributes. For instance, if an element is optional, it definition can be
\begin{alltt}
<xs:element name="\(name\)" type="\(type\)" minOccurs="0" maxOccurs="1"/>
\end{alltt}

An element can have a mixed content (can contain text and other elements at the same time). Such element has to be of a complex type and definition of the type has to contain attribute \texttt{mixed} with value \texttt{true}. The following example demonstrates the mixed content along with the definition of the element type inside \texttt{<xs:element>} schema element.
\begin{alltt}
<xs:element name="\(name\)">
  <xs:complexType mixed="true">
    \(complex type details\)
  </xs:complexType>
</xs:element>
\end{alltt}

The XSD language consists of many more constructs we do not mention such as substitution groups, type extensions, integrity constranints and other.

The book element from the sample XML document in Figure \ref{FIG_a_simple_xml_example} can be described using the XSD as follows.
\begin{alltt}
<xs:element name="book">
  <xs:complexType>
    <xs:sequence>
      <xs:element name="title" type="xs:string"/>
      <xs:element name="authors>
        <xs:complexType>
          <xs:element name="author" type="xs:string" maxOccurs="unbound"/>
        </xs:complexType>
      </xs:element>
      <xs:element name="isbn" type="xs:string" minOccurs="0"/>
    </xs:sequence>
    <xs:attribute name="id" type="xs:string" use="required"/>
  </xs:complexContent>
</xs:element>
\end{alltt}

\section{XPath}
XML Path Language (XPath, version 1.0) \cite{xpath} is a language to select fractions of XML documents using so-called path expressions. XPath considers an XML document as a tree of nodes. It recognizes seven types of nodes: document node, element node, attribute node, text node, namespace node, processing instruction node, and comment node. The root of the tree is the document node representing the entire XML document.

\subsection{Path Expressions}
A path expression is composed of individual path steps separated by \texttt{/} and can be absolute or relative. The absolute path begins with \texttt{/} representing the document node. The relative path needs a non-empty starting node-set to be evaluated. A path step is \\ \texttt{$axis$::$node\_test$[$predicate_1$]$\dots$[$predicate_n$]} \\
where all predicates are optional.

\subsection{Axes}
An axis specifies a node-set relative to the current node. The default axis is \texttt{child} selecting all children nodes of the current node. Another important axes are \texttt{attribute}, selecting attributes of the current node, and \texttt{descendant-or-self} selecting the current node and all its descendants in the tree. The remaining axes are: \texttt{ancestor}, \texttt{ancestor-or-self}, \texttt{descendant}, \texttt{following}, \texttt{following-sibling}, \texttt{namespace}, \texttt{parent}, \texttt{preceding}, \texttt{preceding-sibling}, and \texttt{self}.

\subsection{Node Tests}
A node test identifies node(s) within all nodes selected by an axis. It can be a node type and/or node name. Examples of the node tests follow.
\begin{itemize}
\item \texttt{node()} All nodes selected by an axis.
\item \texttt{text()} Text nodes.
\item \texttt{*} All nodes assigned with their name (element and attributes).
\end{itemize}

\subsection{Abbrevations}
Abbrevations for several most widely used constructs are defined as follows.
\begin{itemize}
\item \texttt{$P$/$identifier$} stands for \texttt{$P$/child::$identifier$} (child is a default axis)
\item \texttt{$P$/@$identifier$} stands for \texttt{$P$/atribute::$identifier$}
\item \texttt{$P$/../$identifier$} stands for \texttt{$P$/parent::*/$identifier$}
\item \texttt{$P$//$identifier$} stands for \texttt{$P$/descendant-or-self::node()/$identifier$} 
\end{itemize}

\subsection{Predicates}
Predicates place additional conditions upon nodes that passed the node test. They can be either relative XPath paths or comparison expressions. The path predicates evaluate to true if they select a not empty set of nodes. Operators in the comparison expressions are \texttt{=}, \texttt{!=}, \texttt{<}, \texttt{>}, \texttt{<=}, \texttt{>=}, \texttt{\&eq;}, \texttt{\&ne;}, \texttt{\&lt;}, \texttt{\&gt;}, \texttt{\&le;}, and \texttt{\&ge;}. Operands are any XPath expressions (paths, literal values, etc).

\subsection{Built-in Functions}
XPath also provides a set of built-in functions. Few examples are \texttt{count($path$)}, returning a number of nodes selected by $path$, \texttt{position()}, returning the position of the current node in the current node-set, and \texttt{sum($path$)}, returning the sum of all nodes selected by $path$.

\subsection{Usage Examples}
Finally, we introduce several examples of XPath expressions, with their description, using the sample XML in Figure \ref{FIG_a_simple_xml_example}.

A path that selects \texttt{author} elements of the book with id $b1$:
\begin{alltt}
/bookstore/book[@id = "b1"]/authors/author
\end{alltt}

An expression returning the number of all books:
\begin{alltt}
count(//book)
\end{alltt}

A path selecting the ISBN element of the book entitled $XQuery$:
\begin{alltt}
/bookstore/book[title = "XQuery"]/isbn
\end{alltt}

A path returning the titles of the books written by more than one author:
\begin{alltt}
/bookstore/book[count(authors/author) > 1]/title/text()
\end{alltt}

\section{XQuery}
An XML Query Language (XQuery, version 1.0) \cite{w3c_xquery} is a language designed to query XML data. It is based on XPath 2.0. XPath 2.0 is an extension of XPath 1.0 (but not entirely compatible), adding ordered sequences, their iterations, set operations, conditional expressions, quantified expressions, and XML Schema types.

In the remainder of this section, basic features of the XQuery languages are briefly described. Many query samples can be found in Chapter \ref{CHAPTER_analysis_of_xquery}.

\subsection{Sequences}
A sequence is an ordered set of items. The result of each XPath 2.0 (and also XQuery 1.0) path is a sequence. An item is either an atomic value or a node. An atomic value is a value of an XML Schema simple type. A node is an instance of one of the node types.

Each node has its identity, data type (simple or complex, according to the XML Schema types), typed value (which can be retrieved by function \texttt{fn:data()}), and string value (which can be retrieved by function \texttt{fn:string()}).

\subsection{FLWOR Expressions}
A basic construct of the XQuery language is $FLWOR$. It is an abbrevation of its five clauses: for, let, where, order by, return. For clause is \\
\texttt{for $var$ in $expr$} \\
Expression $expr$ is evaluated and its result is a sequence. Items of the sequence are iteratively assigned to $var$ variable, which is valid also in the following clauses.

Let clause is \\
\texttt{let $var$ := $expr$} \\
which evaluates $expr$ expression and its result is assigned to $var$ variable.

Where clause is \\
\texttt{where $expr$} \\
Expression $expr$ can (and usually should) contain the variables from the for clause(s) and the remaining clauses are executed for only those tuples of values of the for variables, that the $expr$ evaluates to true.

Order by clause is \\
\texttt{order by $expr$)} \\
and it orders the tuples of values that passed the where clause by the specified criterion.

Return clause is \\
\texttt{return $expr$} \\
The result of the whole $FLWOR$ expression is $expr$. It is constructed using the tuples of the for variable values and the let variables.

\subsection{Conditional Expressions}
XQuery provides common \texttt{if-then-else} conditional expressions with common syntax.
\begin{alltt}
if (\(condition\))
then \(expr1\)
else \(expr2\)
\end{alltt}
If $condition$ expression evaluates to true, $expr1$ is evaluated, else $expr2$ is evaluated. The else branch is optional.

\subsection{Quantified Expressions}
XQuery provides two types of quantified expressions, \texttt{every} and \texttt{some}. Their syntax follows.
\begin{alltt}
every \(var\) in \(expr1\)
satisfies \(expr2\)
\end{alltt}
\begin{alltt}
some \(var\) in \(expr1\)
satisfies \(expr2\)
\end{alltt}

Firstly, $expr1$ is evaluated, and then the result of the quantified expression is true if $expr2$ is true for every (some) item, represented by $var$, of the result of $expr1$ (which is a sequence). 

\subsection{Functions}
XQuery provides a wide set of built-in functions.

And also new functions can be defined using the following syntax.
\begin{alltt}
define function \(name\)(\(parameters\)) as \(type\)
\{
  \(expr\)
\}
\end{alltt}
where $name$ is the function's name, $parameters$ is a list of parameters (with or without specification of their types), and $type$ is a type of a return value, which is the result of $expr$.

%\subsection{Comparison Expressions}
%Comparison expressions are of three types: $value$ comparison expressions, $general$ comparison expressions, and %$node$ comparison expressions. A very brief insight follows. More details can be found in REF\todo{ref, moze to takto ostat?}.

%The $value$ comparison expressions use operators \texttt{lt}, \texttt{gt}, \texttt{le}, \texttt{ge}, %\texttt{eq}, and \texttt{ne} and they compare operands that are or can be converted to atomic values.

%The $general$ comparison expressions use operators \texttt{<}, \texttt{>}, \texttt{<=}, \texttt{>=}, \texttt{=}, %and \texttt{!=} and they can also compare sequences.

%And the $node$ comparison expressions use operators \texttt{is}, \texttt{<<}, and \texttt{>>} and they compare %nodes, their identities and order.

%\section{Other Technologies}

%\begin{description}
%\item[XSLT] e\textbf{X}tensible \textbf{S}tylesheet \textbf{L}anguage \textbf{T}ransformations %\cite{Clark:99:XTV} is a language to transform input XML document(s) into other (not necessarily XML) document(s).
%\end{description}
\chapter{Analysis of Recent Approaches} \label{chapter_analysis_of_recent_approaches}
Approaches of XML schema inference can be divided by several criteria. A basic division is by the language the resulting schema is written in. Commonly used languages are DTD and XML Schema.

According to \cite{Mlynkova:2008:AAX:1494650.1495496}, the type of the inference method can be divided into \emph{heuristic} and \emph{grammar-inferring}. Heuristic methods are based on experience with manual construction of schemas, motivated by real-world usages of XML schema, and their result commonly does not belong to any class of grammar. On the contrary, the grammar-inferring methods are based on theoretical knowledge of automata and their results belong to a particular class of languages. Thus, these methods guarantee specific characteristics of the results. \todo[inline]{Tady by neškodilo dát odkazy na všechny ex. přístupy a říct, že popíšeme detailněji aspoň dva.}

Another important criterion is the type of input data. Most of the approaches process XML documents as the input of the inference process and the documents are supposed to be valid against the resulting schema. Besides approaches exploiting XML data, approaches that utilize other or additional sources may be developed. An approach utilizing XML data along with an obsolete XML schema is described in \cite{Mlynkova:2009:IXS:1862681.1862693}. However, the most significant approaches in terms of this work are those, utilizing operations over XML data. According to our best knowledge at the time of writing, there is just one approach of this category, described in \cite{Necasky:2009:DXK:1529282.1529414}. It utilizes a set of XQuery queries to discover keys and foreign keys\todo{ref do literatury - Ty by tady něškodily všude, u každé té kategorie.}.

\section{Common Caracteristics}
The process of inference commonly used by a significant number of approaches is summarized in \cite{Mlynkova:2008:AAX:1494650.1495496} as the following one: For each occurrence of element $e$ from the input XML documents and its subelements $e_1, e_2, ..., e_k$ a production $e \rightarrow e_1 e_2 ... e_k$ is constructed. The productions form so-called \emph{initial grammar} (IG). For each element type the productions are then merged, simplified and generalized using various methods and criteria. A common approach is so-called \emph{merging state algorithm}, where a \emph{prefix tree automaton} (PTA) is built from the productions of the same element type and the automaton is generalized via merging of its states. Finally, the generalized automaton/grammar is expressed in syntax of the respective XML schema language.
\todo[inline]{Tento odstavec je doslovne prebraty z uvedeneho clanku. Mozem to takto pouzit?}

\section{XTRACT}
The XTRACT \cite{Garofalakis:2000:XSE:342009.335409} system is an example of a \emph{heuristic} \emph{merging state algorithm} creating the result in DTD. Its process of inference consists of three steps:
\begin{enumerate}
\item Generalization - Generates a set of DTD candidates by searching the input for certain patterns and generalising corresponding fragments using regular expressions.
\item Factoring - Groups of generalized candidate DTDs are factorized to a new ones by finding common sub-expressions to make them more concise.
\item Minimum Description Length (MDL) Principle - Composing a near-optimal DTD schema from the set of all generalized candidate DTDs.
\end{enumerate}

\subsection{Generalization}
Purpose of generalization is to create a set of DTD candidates - schemata that cover fractions of the input XML data. In the last step, this set will be used to compose a (sub)optimal result with respect to a trade-off between its preciseness and conciseness. Therefore, it is desirable to create DTD canditates with various degree of these two characteristics.

Generalization is based on replacing fragmets (sequences of subelements of a given element) from the input XML data by regular expressions, thus, using metacharacters like (*, +, ?). To provide a wide set of DTD canditates, each sequence is processed several times using various values of input parameters. Due to the very large number of possible DTD canditates, the authors employ certain real-life motivated heuristics.

For instance, the paper (\cite{Garofalakis:2000:XSE:342009.335409}) introduces the following example: Sequences \texttt{abab} and \texttt{bbbe} are generalized to \texttt{(ab)*}, \texttt{(a|b)*}, \texttt{b*e}.
\todo[inline]{Mozem takto pouzivat prevzate priklady?}

\subsection{Factoring}
Factoring is a process of creating a new DTD canditate from two or more DTD canditates, decresing their summed size without modifications in their semantics. An aim of this step is to decrease a MDL cost of DTD candidates calculated in the MDL step and thus refine the process of construction of the resulting DTD.

An example introduced it the paper (\cite{Garofalakis:2000:XSE:342009.335409}): DTD candidates \texttt{ac}, \texttt{ad}, \texttt{bc} and \texttt{bd} are factored into \texttt{(a|b)(c|d)}.

Alike the generalization step, also in this step the set of possible factored DTD is huge and the authors propose certain heuristics to make the factorization effective.

\subsection{Minimum Description Length (MDL) Principle}
This is an important step trying to create the resulting DTD with the best trade-off between its preciseness and conciseness.  

Paper \cite{Mlynkova:2008:AAX:1494650.1495496} summarizes this step as follows: It expresses the quality of a DTD candidate using two aspects – conciseness and preciseness. Conciseness of a DTD is expressed using the number of bits required to describe the DTD (the smaller, the better). Preciseness of a DTD is expressed using the number of bits required for description of the input data using the DTD. In other words, the more accurately the structure is described, the fewer bits are required. Since the two conditions are contradictory, their balancing brings reasonable and realistic results.

\todo[inline]{Detailnejsie? Priklad?}

\section{Even an Ant Can Create an XSD}
This work, described in \cite{Vosta:2008:EAC:1802514.1802522}, combines several previously proposed approaches including the XTRACT system discussed in the previous section. Its improvements of the process of XML schema inference include:

\begin{itemize}
\item Distinguishing elements with the same name but different semantics.
\item Improvements of algorithms adopted from the previous works.
\item Incorporating inference of an unordered sequence.
\item Creating a result in XML Schema language.
\end{itemize}

\subsection{Clustering of Elements}
This phase clusters elements on a basis of their context and structure. It is done by creating tree structures for each input XML documents, where vertices represent elements and attributes and an edge between two vertices expresses the their parent-child relation. These trees and their subtrees and then compared, using an imposed tree similarity measure, to find elements with the same semantics.

\subsection{Schema Generalization}
For each cluster, a trivial schema is created, which is then generalized to achieve a reasonable result. In a search for the optimal schema, \emph{Ant Colony Optimization (ACO)} heuristic is incorporated. The idea behind the ACO heuristic is that a set of artifical ants is searching a space of possible solutions, each ant given a subspace of the space to find a local suboptimum. An ant 
is performing steps (schema modifications), dying after a predifined number of steps and providing an information - positive feedback - on quality of a found solution. The search is performed in a defined count of interations (or it stops if good enough solution is found) and the positive feedback from one iteration is used to find better results in the following iterations. Every step of an ant represents a modification of a schema, in particular, a merge of states in a corresponding PTA.

One of the improvements of this heuristic is inclusions of a negative feedback after each step of an ant, visible only in a current iteration. Due to this improvement, a larger subspace of solutions is searched.

Another improvement lies in a way how an ant decides for a particular step to perform. To achieve better results, the authors propose a combinations of several verified approaches. A set of all possible steps is created using \emph{k,h-context} and \emph{s,k-string} \todo{citacie} methods. The optimal step is then selected employing the MDL principle \todo{citacia}.

\subsection{Result in XSD}
Unlike the majority of recent approaches, this methods creates its result in XML Schema language (XSD). Comparing to DTD, XSD is a newer language and it provides greater expression possibilities. Authors of this method focused on inferring elements with the same name but different context and the unordered sequences which can be in XSD expressed by \texttt{xs:all} construct. Elements with the same name but different context cannot be expressed in DTD and, however, the unordered sequence can be also expressed in DTD as alternations of ordered sequences, such expression in not practical nor well human-readable.

\section{On Inference of XML Schema with the Knowledge of an Obsolete One}
\todo[inline]{spravne zaradit}
An aim of approach described in \cite{Mlynkova:2009:IXS:1862681.1862693} is to exploit an obsolete XML schema as an additional input information to infer a new schema more efficiently. A XML schema can become obsolete due to changes in a set of XML documents, without capturing these changes in the schema. Thus, the schema becames outdated and according to the paper, this case is quite common.

On input, the method is given:
\begin{itemize}
\item An original XML schema.
\item A set of XML documents. Not all have to be valid against the original schema.
\end{itemize}

And the algorithm consists of two independent steps:

\begin{enumerate}
\item Correction of the input schema.
\item Specialization of the input schema.
\end{enumerate}

In the fist step the input schema is corrected to conform to the whole set of input XML documents. This is done by creating a PTA \todo{ref na def} for each production extracted from the input schema and merging it with a respective production from the initial grammar \todo{ref na def}, involving ACO and MDL heuristics.

In the second step, assuming a correct schema it is wanted to specialize the REs involved in the schema with regard to the XML documents, resulting in more precise and readable schema. Optional substeps are pruning of unused schema fragments, correction of lower and upper bounds of occurrences, transcription of operators to a more restrictive but simplier form, if this transcription preserves the validity, and refactorization to improve readability.

\section{Discovering XML Keys and Foreign Keys in Queries}
This method is described in paper \cite{Necasky:2009:DXK:1529282.1529414}. It tries to improve automatic XML schema inference by discovering keys and foreign keys from a set of XQuery queries. Just the queries are utilized in inference, no XML data are used. The output of this method is a set of keys and foreign keys that can be captured using XML Schema \textbf{key}, \textbf{keyref} and \textbf{unique} constructs.

\subsection{Assumptions and Observations}
To discover keys and foreign keys, the method utilizes element/element joins. Assume a query $Q$ that joins a sequence of elements $S_1$ targeted by a path $P_1$ with a sequence of elements $S_2$ targeted by a path $P_2$ on a condition $L_1 = L_2$.
\todo[inline]{priklad}
The method is based on assumption that each join is done via key/foreign key pair. It means it is supposed that $L_1$ is a key of the elements $S_1$ and $L_2$ is it's respective foreign key or vice versa.

The author describe two possible cases:

\begin{enumerate}
\renewcommand{\theenumi}{(O\arabic{enumi})}
\renewcommand{\labelenumi}{\theenumi}
\item $L_1$ is a key of the elements $S_1$, $L_2$ is a respective foreign key and it itself is not a key of the elements $S_2$.
\item $L_2$ is a key of the elements $S_2$, $L_1$ is a respective foreign key and it cannot be decided whether $L_1$ is a key of the elements $S_1$ or not.
\end{enumerate}

\subsection{Join Patterns and Key Inference}
For a certain join, the decision for one of the cases (O1) and (O2) is made by the form of the join. The query is searched for so-called \emph{join patters}. These are \emph{for join pattern} and \emph{let join pattern} and they are as follows (the first is \emph{for join pattern}, \emph{let join pattern}):

\todo[inline]{spravne kreslit indexy pri pismenkach}
\begin{verbatim}
for $e1 in P1              for $e1 in P1
return                     return
    for $e2 in                 let $e2 :=
        P2[L2 = $e1/L1]            P2[L2 = $e1/L1]
    return CR                  return CR
\end{verbatim}

Each occurrence of a join pattern is classified by application of the following rules it this specific order. The first satisfied ruled is applied. The occurrence is also assigned with a weight determining how sure the method is about the inferred statement.

The pattern occurrence is considered of case (O1) if it is the for join pattern (weight: 1), if aggregation function \texttt{avg}, \texttt{min}, \texttt{max} or \texttt{sum} is applied on a return path (weight: 1) or if aggregation function \texttt{count} is applied on a return path (weight: 0.75). \todo {definovat return path?}

Else, the pattern occurrence is considered of case (O2) and the assigned weight depends on the number of target return paths. If the number is greater than one, the weight is one (weight: 1), else (the number is less or equal one) the weight is one half (weight: 0.5).

\subsection{Summarization of the Results}

The assumption the method is based on may not be fulfilled for every join in a supplied set of queries. A key $K$ may be inferred from some query and processing of another query may result to denial of $K$ as a key. \todo[inline]{priklad?}

Therefore, the authors introduce a scoring function to summarize the positive and the negative statements about keys using the assigned weights. The value of the score expresses the probalitity that a respective key statement is satisfied. Finally, the scores of the inferred keys are normalized to be comparable with each other.

\subsection{Conclusion}
The output of the method is a list of scored keys and for each key a list foreign keys referencing the key. Since the method deals only with the inference of keys, it is not an overall method of the XML schema inference. It is meant to be used in a collaboration with other schema-inferring methods to refine their results.

Since the method is based on intuition of how XQuery constructs are commonly applied in practice, it can be imprecise in certain cases.

\todo[inline]{Tady by se hodila nějaká srovnávací tabulka metod z různách pohledů (metoda, vstupy, výstupy). Dále pak kapitola na téma co dnešní metody neumí a čím se tedy v práci budeme zbývat. Tj. Nějaká návaznost k další kapitole.}
\chapter{Analysis of XQuery}
\todo[inline]{Někde by se mělo říct, kterou verzi uvažujeme, jaké obecně jsou, co je třeba XQuery Core apod.}
This chapter discusses selected constructs of XQuery language and denotes how they could be exploited in the XML schema inference process. It is divided into sections by particular domains of the inference.

Most of sample queries in this chapter are taken from \cite{Walmsley:2007:XQU:1205865} and \cite{Marchiori:07:XQU}, with or without modifications, as these were the main sources for this XQuery analysis.
\todo[inline]{Overit si, ze mozem pouzivat tieto prevzate priklady.}

\section{Structure of XML documents}
Most of queries can be exploited to obtain some information about a structure of respective XML documents. The structure of XML documents captures elements from these documents along with their names, attributes and their organization. What the \todo{Co to znamená? Root element?}outermost elements in these documents are, which elements can be contained within a certain element, whether they are optional or mandatory, etc.

Path expressions without predicates which use only child axis are the simpliest example of such queries. Path expression \texttt{/bib/book/author} indicates that element \texttt{bib} is the outermost element and it contains one or more elements \texttt{book} and these contain one or more elements \texttt{author}.

Additional path expressions \texttt{/bib/book/title}, \texttt{//author/name} indicate that element \texttt{book} also contains element \texttt{title} and element \texttt{author} contains element \texttt{name}. The latter one uses also descendant axis. \todo{A k čemu je tato informace při odvozování?}

Besides elements, attributes can be processed exactly in the same way. Path expression \texttt{/bib/book/@ISBN} indicates that element \texttt{book} has attribute \texttt{ISBN}.

However, statements of this kind are necessary for the XML schema inference, their obtaining from queries is not significant due to the following reason. These statements could be determined directly from the XML documents and it could be easily done in a more convenient way. Queries may not cover the whole relevant content of the documents. For example some element may not be queried at all, therefore they may not be mentioned in any query \todo[inline]{Tady je ale několik případů - chtělo by popsat.
1) Element není zapsaný v dotazu, ale dostanu ho výsledku - např. rok anrození autora, ptám-li se na autory.
2) Element není zapsaný v dotazu, ale procházím ho při vyhodnocování -  např. to //autor
3) Element je v části dokumentu, která se dotazu vlbec netýká.}. In addition, these statements do not express any obligation of occurrence of elements and attributes or possible multiple occurrence of elements \todo{Tak toto je asi taky diskutabilní, např. // by mohlo naznačovat více výskytů.}. We cannot be sure that queries target nodes actually presented in the XML documents. Although query \texttt{/bib/book/author} indicates that element \texttt{author} is contained in element \texttt{book}, the query is valid whether this is true or not. In contrast, even basic methods of XML schema inference that utilize XML documents do not have these inadequacies.

On the other hand, the structural inference could be useful when the entire set of all XML documents is not available and provided XML documents do not cover the structure completely.

\section{Occurrence Count of Elements}
Some XQuery constructs indicate multiplicity of a particular element or limit the element to occur at most once. Consider the following query assuming variable \texttt{\$book1} is bound to a certain \texttt{book} element.

\begin{verbatim}
for $a in $book1/author 
order by $a/last, $a/first
return $a
\end{verbatim}

Apparently, this query expects more than one \texttt{author} element to be a child of the element the variable \texttt{\$book1} is bound to. Otherwise, any sorting would lack a reason. Although we cannot be absolutely sure about it again, assuming common-sense usage of XQuery, it is very likely that \texttt{author} element can occur multiple times as a subelement of element from variable \texttt{\$book1}.

\subsection{Multiple Occurrence}
A similar approach could be applied in many other situations. Another examples are particular usages of function \texttt{count()}, indexation, usage of set operators (\texttt{union}, \texttt{intersect}, \texttt{except}) and usage of function \texttt{one-or-more()}. Sample queries with a respective description follow.

\begin{verbatim}
<section_count>{ count(/book/section) }</section_count>
\end{verbatim}

Function \texttt{count()} returns number of items in a provided sequence. If the sequence is a sequence of elements, the count of these elements will probably not be limited to one. The sample query indicates that the root element \texttt{book} can contain more than one element \texttt{section}.
An exception is usage of function \texttt{count()} in a predicate in expressions where it is used to determine if count of some nodes is greater than zero. Often, this form is used to test presence of a certain node instead of actual counting of its occurrences. In this case, the node could be still limited to occur at most once.

\begin{verbatim}
($s/incision)[2]/instrument
($s/instrument)[position()>=2]
\end{verbatim}

Indexation of nodes and common usage of function \texttt{position()} suggest that the author of such query assumes a sequence of respective nodes.

\begin{verbatim}
one-or-more(/catalog/product[@id = 5]/color)
\end{verbatim}

In this query, number of elements \texttt{color} in element \texttt{product} with attribute \texttt{id} equal to 5 has to be at least one, otherwise an error is raised upon execution. If we assume that this query is written correctly, with common sense and it should not raise the error, we can infer that element \texttt{product} has to contain one or more elements \texttt{color}.

\subsection{Occurrence Limited to One}
Contrary to the multiple occurrence, numerous XQuery constructs limit number of occurrences of an element to at most one or exactly one. Sample queries with description follow.

\begin{verbatim}
/catalog/product[1]/number lt 10
\end{verbatim}

\texttt{lt} is a representant of so called \emph{value comparison operators} \todo{referencia} which operate two sequences of zero or one item. If an operand of a value comparison operator is a sequence of more than one item, then a \emph{type error} \todo{referencia} is raised.

\begin{verbatim}
for $item in //item 
order by $item/num 
return $item
\end{verbatim}

Alike the previous example, an expression in \texttt{order by} clause can be evaluated to at most one item or the \emph{type error} is raised. Therefore, every element \texttt{item} contains zero or one element \texttt{num} but no more.

Another similar examples are arithmetic expressions and functions accepting a sequence of at most one item. Function \texttt{zero-or-one()} will raise the \emph{type error} when supplied with a sequence of more than one item.

Those are constructs indicating limitation to zero or one occurrence. Function \texttt{exactly-one()} works similarly to the function \texttt{zero-or-one()} but accepts only sequences of exactly one item (which are in XQuery equal to this item itself).

\section{Element and Attribute Types}

\subsection{XML Schema Built-in Types}
XML Schema built-in types of elements and attributes can be inferred from the XML documents by analysing their content. Since the count of built-in types is 44 and inheritance is involved, such analysis may be imprecise, especially when a large enough set of XML documents is not available. \todo[inline]{urcite priklad} This is when a set of XQuery queries may come in handy.

If a value of an element or attribute is used in an expression, then this expression could be often exploited to determine a type of the value. Comparing the value to another value of a known type and supplying the value to a function call as an argument of particular type are examples of these expressions as shown below.

\begin{verbatim}
\\event\date = current-date()

\catalog\product\price < 24.5

declare function local:byteFunction($arg as xs:byte) as xs:byte
{...};
\catalog\product\local:byteFunction(@id)
\end{verbatim}
\todo[inline]{overit, ci su priklady OK}

The determined types are \texttt{xs:date} for element \texttt{date}, \texttt{xs:decimal} for element \texttt{price} and \texttt{xs:byte} for attribute \texttt{id}.

Alongside common expressions, there are other XQuery constructs indicating types, such as \emph{type casting}, \emph{constructors} and so called \emph{type declarations}.
\todo[inline]{\^ Nevíme co to je. Chce to vysvětlit. Buď v kapitole 2 nebo tady + odkaz do specifikace.}

Assuming variable \texttt{\$var} bound to a value of some element or attribute, the following two fractions of queries indicate its type to be \emph{xs:integer}.

\begin{verbatim}
$var cast as xs:integer
xs:integer($var)
\end{verbatim}

The latter one is usage of \emph{type casting}, the former one is \emph{type constructor}.

Similarly, the following two queries are examples of \emph{type declarations}. Value of element \texttt{number} is declared to be of type \texttt{xs:integer}.

\begin{verbatim}
every $number as element(*,xs:integer) in //number satisfies ($number > 0) 

declare variable $firstNumber as xs:integer 
:= data(//product/number[1]); 
\end{verbatim}

Although determining built-in types of elements and attributes in simple queries may seem quite straightforward, reliable method would face some issues. More complex queries are commonly branched by \emph{if-then-else} construct. These branches can contain several contradictionary fragments from the view of type inference. Therefore, the reliable method would take the query structure into consideration and further research would be needed.

\subsection{Enumeration}
In several cases, XQuery \emph{if-then-else} construct may be used to branch the execution of query by all possible values of a certain variable. This is equivalent of \emph{switch} construct from other programming languages. When the control variable is bound to some element or attribute, type of this node can be inferred as enumeration and its individual values can be determined as well.

\todo[inline]{Tady by to chtělo příklad. nebo to budeme řešit dál v navrženém algoritmu a toto je pouze diskuse případů?}

\begin{verbatim}
Query
let $cat := doc("catalog.xml")/catalog 
for $dept in distinct-values($cat/product/@dept) 
return <li>Department: {if ($dept = "ACC") 
                        then "Accessories" 
                        else if ($dept = "MEN") 
                             then "Menswear" 
                             else if ($dept = "WMN") 
                                  then "Womens" 
                                  else () 
  } ({$dept})</li> 

Results 
<li>Department: Womens (WMN)</li> 
<li>Department: Accessories (ACC)</li> 
<li>Department: Menswear (MEN)</li> 
\end{verbatim}

Finding such patterns in queries could be useful in combination with analysis of respective XML data. XML data could help to confirm or disprove this assumption based on the query analysis or vice-versa.

\section{Keys}

\subsection{Approach from [TODO citace]}
Paper \todo{link} introduces a method of discovering keys and foreign keys by investigation of joins in queries. The basis of this discovery is a search for particular forms of joins, so called \emph{join patterns} \todo{uviest ich sem?}, but the joins are processed only if they are found in a particular syntactic form. Therefore joins with the same semantics written in different syntax are not taken into consideration.

The following examples are queries that could be processed by a similar method but the actual method will not use them.

\begin{verbatim}
<result>
  {
    for $u in doc("users.xml")//user_tuple
    for $i in doc("items.xml")//item_tuple
    where $u/rating > "C" 
       and $i/reserve_price > 1000 
       and $i/offered_by = $u/userid
    return
        <warning>
            { $u/name }
            { $u/rating }
            { $i/description }
            { $i/reserve_price }
        </warning>
  }
</result>
\end{verbatim}

Clause \texttt{where} is used for the join condition \texttt{\$i/offered\_by = \$u/userid instead} of the join condition in predicate of the second \texttt{for} expression.

\begin{verbatim}
<result>
  {
    for $i in doc("items.xml")//item_tuple
    where empty(doc("bids.xml")//bid_tuple[itemno = $i/itemno])
    return
        <no_bid_item>
            { $i/itemno }
            { $i/description }
        </no_bid_item>
  }
</result>
\end{verbatim}

This query illustrates a join where the value of one of the joined elements is not required, thus the knowledge of its existence is sufficient. Therefore the query does not have to contain the second \texttt{for} or \texttt{let} keyword and its expression can be moved to \texttt{where} clause.

\begin{verbatim}
for $item in doc("order.xml")//item, 
    $product in doc("catalog.xml")//product, 
    $price in doc("prices.xml")//prices/priceList/prod 
where $item/@num = $product/number and $product/number = $price/@num 
return <item num="{$item/@num}" 
         name="{$product/name}" 
         price="{$price/price}"/>
\end{verbatim}

A common three-way join.

\subsection{Join of Recursive Data}
By the term ``recursive data'' are meant XML elements that somehow reference items from the same set. \todo[inline]{Tohle je velmi nepřesné. Tady asi myslíte reference, ale rekurzivní element je obecně element, kteý obsahuje sám sebe - např. kapitola má kapitoly.} Example of a query that operates upon such data follows.

\begin{verbatim}
declare function local:one_level($p as element()) as element()
{
  <part partid="{ $p/@partid }" name="{ $p/@name }" >
    {
      for $s in doc("partlist.xml")//part
      where $s/@partof = $p/@partid
      return local:one_level($s)
    }
  </part>
};

<parttree>
  {
    for $p in doc("partlist.xml")//part[empty(@partof)]
    return local:one_level($p)
  }
</parttree>
\end{verbatim}

Apparently, element \texttt{part} can contain attributes \texttt{partid} and \texttt{partof}. Elements with unspecified attribute \texttt{partof} are at the top of the recursive hierarchy while, in elements with this attribute specified it makes a reference to an element with attribute \texttt{partid} of the same value. \todo{Tahle věta je nějaká divná.}

By using a similar approach as described in the paper \todo{ref} attribute \texttt{partid} can be marked as a key and attribute \texttt{partof} as its foreign key.

\subsection{Negative Statements about Uniqueness}
In many cases a statement refusing uniqueness of element or attribute values could be inferred. Such statements may be helpful in combination with other methods in its process of making a decision whether a particular element or attribute is unique or not.

Basic representant is a common \emph{FLWOR} query.

\begin{verbatim}
<bib>
  {
    for $b in doc("http://bstore.example.com/bib.xml")//book
    where $b/publisher = "abcde" and $b/@year > 2000
    order by $b/title
    return
        <book>
            { $b/@year }
            { $b/title }
        </book>
  }
</bib>
\end{verbatim}

Usage of \texttt{for} construct indicates that a sequence of \texttt{book} elements which satisfy the condition in \emph{where} clause is expected. It is a condition composed of two single conditions joined by \emph{and} logical operator. Therefore, in order to satisfy the whole condition, both single conditions must be satisfied as well. Thus, it is expected that several elements satisfy each of the single conditions.

The first of them is test of equality of book's subelement \texttt{publisher} to a string literal. Based on the expectation, element \texttt{publisher} cannot be unique. However, the second condition is greater-than comparison of \texttt{year} attribute to an integer literal, it cannot be inferred whether this attribute is unique or not. The reason is that even if it was unique, there still might be more than one \texttt{book} element meeting this condition. Also, any statement, positive nor negative, cannot be inferred about \texttt{title} subelement.

Other simple examples are passing a result of basic path expression to a call of \texttt{distinct-values()} function and usage of \emph{aggregation} functions.

\begin{verbatim}
<results>
  {
    let $doc := doc("prices.xml")
    for $t in distinct-values($doc//book/title)
    let $p := $doc//book[title = $t]/price
    return
      <minprice title="{ $t }">
        <price>{ min($p) }</price>
      </minprice>
  }
</results>
\end{verbatim}

According to the use of \texttt{distinct-values()} function, it can be easily seen that the author of this query assumes possible occurrence of more than one \texttt{book} elements with the same value of their \texttt{title} subelement. Thus, \texttt{title} element is not unique.

Alike, variable \texttt{\$p} is bound to a price of each title and then passed to \texttt{min()} function call. That indicates that one book title is supposed to have several prices, however, it cannot be said if a certain occurrence of element \emph{book} can contain more than one \texttt{price} subelement.

Also, many other types of queries could be exploited to obtain negative uniqueness statements like occurrence of element or attribute in \texttt{stable order by} clause or selection of an element set based on a value of particular attribute and consecutive treatment of this set as a sequence of elements.

\begin{verbatim}
let $prods := doc("catalog.xml")//product 
for $prod in $prods 
where $prod << $prods[@dept = $prod/@dept][last( )] 
return $prod 
\end{verbatim}

\subsection{Uniqueness}
Contrary to non-uniqueness discovery, some XQuery constructs can indicate uniqueness of elements or attributes; however, it seems to be more difficult. One of the approaches of uniqueness discovery could be an investigation of what the query does return. Consider the following query.

\begin{verbatim}
for $product in /catalog/product
let $number := $product/number
return <prod xml:id="{concat('prod', number)}"/>
\end{verbatim}

For each \texttt{product} element, new element \texttt{prod} with attribute \texttt{xml:id} is created. Since attribute \texttt{xml:id} is supposed to be unique and there is direct transformation of the values of \texttt{number} elements to the values of attribute \texttt{xml:id}, it is very likely that element \texttt{number} is unique in the source data.

\section{Other Constructs}

\subsection{XML Schema \texttt{xs:sequence} and \texttt{xs:all} constructs}
If an order of appearance of some element set, for example subelements of a certain element, is important, this can be expressed by XML Schema construct \texttt{xs:sequence}. On the other hand, \texttt{xs:all} is involved when the elements may occur in any order.

If there is a large enough set of XML documents available, it should be possible to correctly detect where to use \texttt{xs:sequence} and \texttt{xs:all} constructs. However, if the inference is made using a smaller set of XML documents, it may happen that every occurrence of some element set is in the same order but the input data are not representative enough to be sure. Comparison of element order in queries can help to decide for the use of \texttt{xs:sequence}.

Every two elements in an XML document that are not siblings with the same name and none of them is the root element have the nearest common ancestor determining their relative order. If there is a lack of evidence to choose between the two constructs in the ancestor, then available queries can be searched for an order comparison of these two elements like in the following query.

\begin{verbatim}
let $i := //incision)[2]
for $a in //action)[. >> $i]
return $a//instrument
\end{verbatim}

Apparently, the relative order of elements \texttt{incision} and \texttt{action} is important, because the query composes the result using only those \texttt{action} elements that succeed the second \texttt{incision} element. Therefore, in a respective type definition part of their nearest common ancestor \texttt{xs:sequence} should be used in favour of \texttt{xs:all}.

\subsection{Intermediate XML structure}
Intermediate XML structure are XML data that are not read from input XML documents nor created as an output of queries but they are somehow used by the queries. Objective of this section is intermediate XML structure created directly in the queries. It may serve to various purposes and it can be created in various ways, hard-coded in queries, computed from the input XML documents to simplify their structure. The former is demonstrated in the following example.

\begin{verbatim}
let $deptNames := <deptNames> 
                    <dept code="ACC" name="Accessories"/> 
                    <dept code="MEN" name="Menswear"/> 
                    <dept code="WMN" name="Womens"/> 
                  </deptNames> 
let $catalog := doc("catalog.xml")/catalog 
for $dept in distinct-values($catalog/product/@dept) 
return <li>Department:
         {data($deptNames/dept[@code = $dept]/@name)} ({$dept})
       </li> 
\end{verbatim}

Intermediate XML structure in this query can be used to identify possible values of attribute \texttt{dept} in \texttt{product} element and therefore determine its type as enumeration.

This is a very simple example of intermediate structure. Since there are admittedly many means of utilization of intermediate structure, further research would be needed.
\chapter{Pre-creation of Algorithm}
According to the analysis in the previous chapter, there is quite a wide range of possible utilizations of XQuery queries. Besides analysis of what information can be extracted from queries, it is needed to devise how the queries will be processed. This chapter discusses some questions and issues that emerged in an early phase of the algorithm fabrication.

\section{Input Data}
The first important question is what is the input of the algorithm. A basic query utilization can be achieved by analysis of queries without any other input data. The analysis of XQuery in the previous chapter discusses mostly XQuery constructs which can be utilized without respective XML data, for example the inference of built-it types. This independence is also the main advantage of this approach, if there are no XML data available, this approaches can be still used. 

A more complex method can utilize queries along with the respective XML data. As discussed in the previous chapter, an element and attribute structure can be inferred from the XML data in a more convenient way. Also, the XML data can be used to verify information inferred from the queries or vice-versa. For example, utilizing the queries, some attribute is considered a key of its element. But in the data there are elements with the same value of this attribute, and thus, it cannot be the key. Vice-versa, we have a notion that the attribute might be the key but we are not sure about that. Analysing of the data and finding that values of the attribute are unique can increase our confidence. 

Another step can be evaluation of the queries using the XML data and consecutive analysis of the results. And even the process of evaluation itself can be analysed to obtain some useful information. For instance, these are partial results of evaluating of expressions (elements selected by each path expression, real arguments in function calls, etc).

\section{Forms of Query Precessing}
Another important question is how the queries can be processed. Will they be just searched for certain patterns like it is performed in method \cite{Necasky:2009:DXK:1529282.1529414} or will they be processed in a more sophisticated way? That can mean incorporating lexical and syntax analyses, known from creation of compilers, or even a form of an analysis of semantics \cite{compilers}.

The result of lexical and syntax analyses can be a kind of so-called syntax tree \cite{compilers}. It is a structure representing a word according to a formal grammar of a language. In our case, the language is XQuery, its grammar is defined in \cite{w3c_xquery} and every query is a word of the XQuery language. Leaves of the tree represent terminals of the grammar while inner nodes represent non-terminals. \todo{priklad nejakeho stromu?} From the point of view of this work, the syntax tree can be perceived as a preprocessed form of a query keeping its complete meaning and making its further processing more convenient. For instance, the tree can simplify a search for FLWOR statements. It is transitioned and nodes representing FLWORs are found. Then each subtree determined by one of the found nodes represents a FLWOR statement and it can be analysed further.

\todo[inline]{Q: definovat grammar, tree, language? A: Pokud bude použito dále, tak ano.}

The syntax tree can be also extended by additional information. An example is a static analysis of expression types. Types of literal expressions are defined, functions have return types, path expressions can return nodes, etc. Types of complex expressions can be determined applying the rules defined in \cite{w3c_xquery}. The inferred expression types can be helpful for example in the analysis of built-in types of nodes as discussed in the previous chapter. 

The following text is an example of inference of a more complex query processing. Consider the following part of a query.
\begin{verbatim}
declare function local:getB($id as xs:string) as element() {
  //A[@id = $id]/B
}
... local:getBs("id") > 10 ...
\end{verbatim}
The query consists of a function declaration and an arithmetic comparison. The comparison compares a result of the function call and literal value \texttt{10}. For the type of \texttt{10} is \texttt{xs:integer} the type of the function call has to be convertible to \texttt{xs:double}. Thus, it has to be a numeric type. The function returns a path expression typed as \texttt{element()}. That means the function returns one element. In the path expression, the argument \texttt{\$id} can be substituted by the real value specified in the call. Thus, the return expression is \texttt{//A[@id = "id"]/B}. Therefore, we can infer that element \texttt{B} in element \texttt{A} with attribute \texttt{id} equaling string value \texttt{"id"} is of some numeric type. And, since the function is parametrized, there is a notion that this statement may be correct also for other elements \texttt{A} and \texttt{B}.

While simplifier approaches of the query processing such as the pattern finding limit possibilities of the query utilization, a more complex processing of queries provides a better starting point for consecutive analyses and also for further refinements and additions. Therefore, we decided to incorporate query processing using the lexical and syntax analyses.

\section{Inference of XML Structure}
The question of inference of XML structure from queries is partially discussed in the previous chapter. We are able to infer XML structure from queries without their evaluation, but in a limited way. This inference is based on an analysis of path expressions. Its limitation involves the following issues. When we infer some subelements of a certain element we often cannot be sure about their number of occurrences. Also, we cannot be sure if every occurrence of a certain element contains the subelements and we even do not know if at least one occurrence of the element contains them. Thus, the inferred statement is more likely an indication than a fact about the structure.

Since the inference of XML structure utilizing only queries is not clear, we need the XML data, if we want to infer the structure more precisely. And, if we have the XML data, we can infer the structure directly from them using an existing approach and utilize queries to refine its result.

\section{Extension of an Existing Approach}
The existing approaches of XML schema inference deal mainly with inference of XML structure. Hence, the extension of an existing approach will be a kind of an independent addition instead of modification and refinement of its core algorithm.

The existing approaches take the XML data on their input. Therefore, the basic idea is that the input will be extended also for XQuery queries and the algorithm will consist of three phases. The first phase will be taken from an existing approach and it will process the XML data to infer the XML structure. The second phase will process the XQuery queries and it will infer statements that can be inferred independently of the XML data. The third phase will merge the statements inferred in the second phase into the resulting schema. This phase may also infer additional statements from both the XML data and the queries or it may try to verify the statements from the second phase with respect to the XML data.

A more advanced method can exploit queries to refine a core algorithm (e.g. merging of PTA) of an existing approach. The approach described in \cite{Vosta:2008:EAC:1802514.1802522} distinguishes elements with the same name, but a different content model and context. Some information from queries may help to improve this algorithm. For example, consider an XML representation of company data containing names of employees, costumers, and products represented by element \texttt{name}. The names of employees and costumers consists of two subelements for first name and surname. The names of products are atomic strings. During analysis of available queries, we may find that elements representing the names of employees and elements representing the names of costumers are processed in the same way (e.g. they are mixed in one sequence), while elements representing the names of products are processed separately. This suggests that the elements representing the names of employees and costumers have the same semantic and the same content model which is different from the content model of the elements representing the names of products.
\chapter{Proposed Algorithm}
\todo[inline]{Nejaky uvod a zkratke hrnutie}
\todo[inline]{Lexikalna a syntakticka analyza z diplomky J. Schejbala}
\todo[inline]{Staticka analyza typov vyrazov}
\todo[inline]{Predpoklad korektnosti dotazov}

\section{Step 1: Construction of a Syntax Tree}
The first step of the algorithm involves lexical and syntactic analyses known from compilers and produces a so-called syntax tree. The analyses are taken from Ji\v r\'{i} Schejbal's master thesis \todo{ref}. Since they are not directly related to the inference, and thus, they are not directly related to the topic of this thesis, we will not describe them. Nevertheless, they provides us with a helpful processing of XQuery queries and we can focus on the inference.

The syntactic analysis needs to be slightly modified to suit our case. It writes its result into a file in an XML representation. Instead, we need to keep the result in memory and pass it to consecutive steps of our algorithm. Though this requirement concerns modifications of implementation, the core of the processing remains untouched. Therefore, we also do not describe these modifications.

\todo[inline]{Toto vše je OK, ale asi budeme potřebovat vědět, co nám z té analýzy leze za informace se kterými dál pracujeme. Čili vstupy našeho algoritmu je potřeba popsat přesně, ideálně formálně, tj. je to asi nějaký strom T= (V, E), kde V je... atd. Jeslti už je to pak implementováno jako objekty v paměti nebo XML data je celkem jedno a bude popsáno v kapitola Implementation.}

\subsection{Definition of the Syntax Tree}
\begin{definition}
Syntax tree of a XQuery query $Q$ is a tuple $T = (V, E, o)$ where $V$ is a set of nodes representing XQuery language constructs (TODO\todo{kompletny zoznam}), $E$ is a set of edges - couples $(a,b):a,b \in V$ and $o:V \rightarrow \mathbb{N}$ is an ordering function and $T$ satisfies the following conditions:
\begin{itemize}
\item $(V,E)$ is a tree.
\item The root node is $ModuleNode$ representing the entire query Q.
\item $(a,b)$ is an element of $E$ if and only if $b$ is a direct component of construct $a$ in $Q$.
\item For every node $a \in V$ the ordering function $o$ is defined as follows \\* $o(a) = \left\{
	\begin{array}{ll}
		1  & \mbox{if }a\mbox{ is the root node} \\
		seq(a) & \mbox{else}
	\end{array}
\right.$ \\* where $seq(a)$ is a sequential number of appearence of construct $a$ in $Q$.
\end{itemize}
\end{definition}

\subsection{Syntax Tree Nodes}
\todo[inline]{Dat do appendixu?}
Nodes of the syntax are organized in an is-a hierarchical structure, commonly used in the object oriented programming languages, where a simple node can be of several types. Nodes with emphasised names express a common type of a group of nodes - their descendants, and these type nodes cannot be directly used in the syntax tree. The remaining nodes represent particular constructs of the XQuery language and their instances can be used as nodes of the syntax tree. Each of the instances additionaly contains details about the particular form of the respective XQuery construct used in the query.

\todo[inline]{V ExprNode s ExprHolderNode si pamatame aj typ vyrazu}

For example, an instance of the syntax tree cannot directly contain nodes $Node$ and $ExprNode$, but it can contain instances of nodes $AttributeNode$ and $LiteralNode$. Regarding the types, an instance of $LiteralNode$ node is of three types: $LiteralNode$, $ExprNode$ and $Node$. It carries information on a type of the literal and its value, for example $xs:integer$ of value $5$.

Nodes can be classified into three groups: inner nodes, leaf nodes and nodes that can be both inner and leaf. Inner nodes (marked $I$) stand for XQuery constructs that are composed of other constructs and can be further divided. An example of such node is $FLWORExprNode$ which is composed of nodes $TupleStreamNode$, $WhereClauseNode$, $OrderByClauseNode$ and $ReturnClauseNode$. Leaf nodes (marked $L$) represent elements of XQuery language that cannot be further divided. For example, $LiteralNode$. Node $FunctionCallNode$ is an example of the third group (marked $IL$). Function call of a function without arguments is represented by a leaf while function call with arguments is represented by an inner node and its subnodes are nodes representing those arguments.

\begin{description}
\item \textbf{Node}
\begin{description}
\item AttributeNode (I)
\item AttrListNode
\item AxisNode (I)
\item CaseClauseNode (I)
\item CaseClausesNode (I)
\item CDataSectionNode
\item ContentNode
\item DefaultCaseNode (I)
\item EntityRefNode
\item \textbf{ExprHolderNode}
  \begin{description}
  \item BindingSequenceNode (I)
  \item DefaultValueNode (I)
  \item ElseExpressionNode (I)
  \item OperandExpressionNode (I)
  \item ReturnClauseNode (I)
  \item TestExpressionNode (I)
  \item ThenExpressionNode (I)
  \item WhereClauseNode (I)
  \end{description}
\item \textbf{ExprNode}
  \begin{description}
  \item CommaOperatorNode (I)
  \item ConstructorNode (I)
  \item ContextItemExprNode
  \item EmptySequenceNode
  \item ExtensionExprNode (I)
  \item FLWORExprNode (I)
  \item FunctionCallNode
  \item IfExprNode (I)
  \item LiteralNode
  \item OperatorNode (I)
  \item OrderedExprNode (I)
  \item PathExprNode (I)
  \item QuantifiedExprNode (I)
  \item TypeswitchExprNode (I)
  \item UnorderedExprNode (I)
  \item ValidateExprNode (I)
  \item VarRefNode
  \end{description}
\item FunctionBodyNode
\item CharRefNode
\item InClausesNode
\item \textbf{ItemTypeNode}
  \begin{description}
  \item AnyItemNode
  \item AtomicTypeNode
  \item KindTestNode
  \item NameTestNode
  \end{description}
\item LocationHintNode
\item LocationHintsNode
\item \textbf{ModuleChildNode}
  \begin{description}
  \item ModuleDeclNode
  \item PrologNode
  \item QueryBodyNode
  \end{description}
\item ModuleNode
\item NameNode
\item OrderByClauseNode
\item OrderSpecNode
\item ParamListNode
\item ParamNode
\item PITargetNode
\item PragmaListNode
\item PragmaNode
\item PredicateListNode
\item \textbf{PrologChildNode}
  \begin{description}
  \item BaseURIDeclNode
  \item BoundarySpaceDeclNode
  \item ConstructionDeclNode
  \item CopyNamespacesDeclNode
  \item DefaultCollationDeclNode
  \item DefaultNamespaceDeclNode
  \item EmptyOrderDeclNode
  \item FunctionDeclNode
  \item \textbf{ImportNode}
    \begin{description}
    \item ModuleImportNode
    \item SchemaImportNode
    \end{description}
  \item NamespaceDeclNode
  \item OptionDeclNode
  \item OrderingModeDeclNode
  \item VarDeclNode
  \end{description}
\item StepExprNode
  \begin{description}
  \item SelfOrDescendantStepNode
  \end{description}
\item StringNode
\item TupleStreamNode
\item TypeNode
\item \textbf{VariableBindingNode}
  \begin{description}
  \item ForClauseNode
  \item InClauseNode
  \item LetClauseNode
  \end{description}
\item VarValueNode
\end{description}
\end{description}

Some pseudo code algorithms in the following sections need to determine the type of a node instance. For this purpose, we define function $isA(nodeInstance, nodeType)$. The first argument is an instance of any node and the second is a node type itself. The function returns boolean value $true$ if the instance is of the specified node type, according to the described principle. Otherwise, it returns $false$.

For instance, assuming variable $literal$ is an instance of node $LiteralNode$, calls $isA(literal, LiteralNode)$, $isA(literal, ExprNode)$ return $true$, while call $isA(literal, ContentNode)$ returns $false$.

\subsection{Description of the Syntax Tree}
Since we modified only the implementation of the output form, the design of the syntax tree remains consequent to the entire design of the analyses. Each node represents a logical component of a query. 

Inner nodes are: $AttributeNode$, $AxisNode$, $CaseClauseNode$, $CaseClausesNode$, $DefaultCaseNode$, $BindingSequenceNode$, $DefaultValueNode$, $ElseExpressionNode$, $OperandExpressionNode$, $ReturnClauseNode$, $TestExpressionNode$, $ThenExpressionNode$, $WhereClauseNode$, $CommaOperatorNode$, $ConstructorNode$, $ExtensionExprNode$, $FLWORExprNode$, $IfExprNode$, $OperatorNode$, $OrderedExprNode$, $PathExprNode$, $QuantifiedExprNode$, $TypeswitchExprNode$, $UnorderedExprNode$, $ValidateExprNode$, $InClausesNode$, $QueryBodyNode$, $ModuleNode$, $OrderByClauseNode$, $OrderSpecNode$, $ParamListNode$, $ParamNode$, $PragmaListNode$, $PredicateListNode$, $FunctionDeclNode$, $ModuleImportNode$, $SchemaImportNode$, $VarDeclNode$, $StepExprNode$, $TupleStreamNode$, $ForClauseNode$, $InClauseNode$, $LetClauseNode$.

Nodes that can be both inner and leaf are: $AttrListNode$, $ContentNode$, $FunctionCallNode$, $FunctionBodyNode$, $KindTestNode$, $LocationHintsNode$, $NameNode$, $PrologNode$, $PITargetNode$, $VarValueNode$.

Leaf nodes are: $CDataSectionNode$, $EntityRefNode$, $ContextItemExprNode$, $EmptySequenceNode$, $VarRefNode$, $PragmaNode$, $BaseURIDeclNode$, $BoundarySpaceDeclNode$, $ConstructionDeclNode$, $CopyNamespacesDeclNode$, $DefaultCollationDeclNode$, $DefaultNamespaceDeclNode$, $EmptyOrderDeclNode$, $LocationHintNode$, $NamespaceDeclNode$, $OptionDeclNode$, $OrderingModeDeclNode$, $SelfOrDescendantStepNode$, $StringNode$, $TypeNode$.

\todo{zistit, ci konstanty nahodou nie su jedine. NIE: CDataSectionNode, ContentNode(ale nemusi byt list), EntityRefNode, ContextItemExprNode, EmptySequenceNode, FunctionCallNode(bez parametrov), VarRefNode, FunctionBodyNode(prazdna funkcia), CharRefNode, AnyItemNode, AtomicTypeNode, KindTestNode(nemusi), NameTestNode, LocationHintNode, ModuleDeclNode, PrologNode, NameNode(nemusi), PITargetNode(nemusi), PragmaNode, BaseURIDeclNode, BoundarySpaceDeclNode, ConstructionDeclNode, CopyNamespacesDeclNode, DefaultCollationDeclNode, DefaultNamespaceDeclNode, EmptyOrderDeclNode, NamespaceDeclNode, OptionDeclNode, OrderingModeDeclNode, SelfOrDescendantStepNode, StringNode, TypeNode, VarValueNode}

\todo[inline]{\^ Toto bude potřeba popsat přesně, ne jen příkladem. Klidně převezměnte z diplomky JS.}

In detail, the syntax tree is a n-ary tree representing syntactic structure of a given query. Every node represents a certain XQuery construct and that is a type of the node. Some types can be divided into more specific subtypes. This is shown \todo{obrazok typov uzlov}.

\todo[inline]{\^ Toto je skoro stejný odstavec ne?}

An important characteric of the syntax tree is related to a definition of local variables and their scope. The representation of a definition of a local variable is $VariableBinding$ node. Nodes of this type contain only two subnodes; a type of the variable and a binding expression. The node and its subtree representing the type do not contain any expressions. The binding expression node is an expression node but the represented expression cannot use the variable which is just being bound. Therefore the scope of the new variable is not the subtree of the $VariableBinding$ node. Assuming \todo{Toto je velmi high-level popis na to, jak málo o syntaktickém stromě a dalších položkách víme.} the post-order numbering of the nodes, the new variable may be valid in some nodes with a higher post-order sequence number. Which \todo{Toto není anglický slovosled.} nodes that are, it depends on the type of the ancestor node(s).

For example, $FLWORExpr$ node contains four subnodes; $TupleStream$, $WhereClause$, $OrderByClause$ and $ReturnClause$. $TupleStream$ node contains a list of $VariableBinding$ node which define variables valid in all other three subnodes of the $FLWORExpr$ node.

\todo[inline]{priklad stromu}

\section{Step 2: Static Analysis of Expression Types}
In the second step, we statically (i.e. without evaluation of the queries) determine types of expressions in the syntax tree. Information on the types of expressions can be used by consecutive steps of the algorithm. The consecutive steps will not use the determined types of all expressions, however this step may be useful in a future extending.

\todo[inline]{Tady z těch try a probably je jasně vidět, že netušíte co vlastně přesně algoritmus bude dělat. Nejdříve si to rozmyslete, pak pište text. Rozhodně by se takové výrazy neměly v DP dobjevovat, působít o velmi nedůvěryhodně. Měl by to být strohý popis toho, co bylo uděláno.}

We want to capture the following types.

\begin{itemize}
\item XML Schema built-in atomic types \todo{link do appendixu}.
\item Type $xs:untyped$. \todo{to je co?}
\item Type representing any item. \todo{je potrebne?}
\item Type representing an element, attribute, item, node, text node, comment, processing instruction, document node.
\item Type representing a set of elements or attributes selected by a certain path expression. The path expression is included in this type. Let this type be identified as $pathType$.
\end{itemize}

\todo[inline]{Q: Akeho typu je napr path expression /root/(elementA * 2)? A: Datový typ je daný operátorem. Hodnota elementu se přetypuje.}

To capture sequences, we assign each type with its cardinality as proposed in \todo{ref na Schajbalovu diplomku}. Every type can be perceived as a sequence. A type representing one value or one item can be perceived as a sequence of exactly one item. The cardinality expresses one of the following five sequence types.

\begin{itemize}
\item An empty sequence.
\item A sequence of exactly one item.
\item A sequence containing zero or one item (modifier $?$).
\item A sequence containing zero or more items (modifier $*$).
\item A sequence containing one or more items (modifier $+$).
\end{itemize}

If we cannot determine the type of an expression, we express that using $null$. An example of such case is path expression $/rootElement/(elementA * 2)$. \todo{Nejak dopisat}

The analysis of expression types can be divided into three substeps. Determination of return types of functions, determination of types of global variables, and finally determination of types of expressions.

\todo[inline]{Este by sa dala urobit analyza volania funkcii, do ktorych su predavane nejake elementy/atributy o ktorych vieme ich path expression. Vysledkom tejto funkcie by totiz zas mohli byt elementy/atributy o ktorych vieme (pozmenenu) path expression a nie len nezname elementy/atributy.
A: Přesně tak. A jak říkám - není třeba vše co navrhnete implementovat.}

\subsection{Determination of Function Return Types}
Determination of return types of functions is needed because function calls can appear in expressions. A return type of a function can be determined at the moment the analysis of expressions encounters a call of the function, however, it involves multiple transitions of the syntax tree in a search for a definition of a particular function.

Instead, the syntax tree can be searched just once, before the analysis of expressions, and return types of all functions found are stored.

A simple algorithm is presented in Algorithm \ref{ALG_determination_of_function_return_types} \todo{Q: uvadzat aj taketo jednoduche funkcie? A: Obecně by mělo být popsáno vše. Ale pokud se nějaká jednoduchá věc např. opakuje, tak stačí napsat, že je algoritmus podobný jako v předchozím případě nebo tak něco.}. It uses \texttt{getFunctionDeclarationNodes()} function, which returns a list of all function declaration nodes in the syntax three. Actually, it does not have to search the whole syntax tree as the function declaration nodes can be present only in the query prolog section. 

For simplicity, we will focus on locally defined functions. We will not determine types of functions defined in other modules since the principle is similar but it is somewhat more difficult to implement \todo{Q: je to ok? A: No spíš napište něco jako without loss of generality.}. Return types of built-in functions are fixed and, thus, there is no need to analyze them.

\begin{algorithm}
\caption{Determination of Function Return Types}
\label{ALG_determination_of_function_return_types}
\begin{algorithmic}[1]
\ENSURE An associative array of function names with their types.

\STATE $functionTypes$ = empty associative array
\FORALL{$functionDeclarationNode \in getFunctionDeclarationNodes()$}
    \STATE $functionName = functionDeclarationNode.getFunctionName()$
    \STATE $functionTypes[functionName] = functionDeclarationNode.getReturnType()$
\ENDFOR
\RETURN $functionTypes$
\end{algorithmic}
\end{algorithm}

For the rest of the thesis, we assume function \todo{Většinou se na začátek sekce s algoritmem dává seznam používaných  funkcí a jejich význam.} \texttt{getFunctionReturnType()} which takes a function name and returns the return type of the function if the function is either built-in or it is recorded in the associative array. Otherwise, \texttt{null} is returned.

Later phases of the algorithm may need to process the function declarations. Therefore, we also store references to the entire function declaration nodes.

\subsection{Determination of Global Variable Types}
A similar approach as in the case of functions can be applied to determine types of global variables. Alike the functions, the global variables are defined in the prolog section. A type of a variable can be explicitly specified in its definition \todo{je to tak? priklad}. If it is not, it may often be deducible from the binding expression. Again, we do not analyse external variables for the same reason as in case of external functions.

\begin{algorithm}
\caption{Function determineExpressionType}
\label{ALG_func_determineExpressionType}
\begin{algorithmic}[1]
\REQUIRE{\ \\
    $exprNode$: An expression node in the syntax tree. \\
    $contextVarTypes$: Types of local variables valid in the current subtree.
}
\ENSURE The expression type.

\STATE $type = null$
\IF{$isA(exprNode, LiteralNode)$}
     \STATE $type =$ ((LiteralNode)exprNode).type, exactly-once cardinality
\ELSIF{$isA(exprNode, FunctionCallNode)$}
     \STATE $type = getFunctionReturnType(((FunctionCallNode)exprNode).functionName)$
\ELSIF{$isA(exprNode, VarRefNode)$}
     \STATE $type = getVariableType(contextVarTypes, ((VarRefNode)exprNode).variableName)$
\ELSIF{$isA(exprNode, PathExprNode)$}
     \STATE $path = ((PathExprNode)exprNode).path$
     \IF{$path$ selects an element set}
         \STATE $type =$ element type, zero-or-more cardianlity, $path$ included
     \ELSIF{$path$ selects an attribute set}
         \STATE $type =$ attribute type, zero-or-more cardianlity, $path$ included
     \ELSE
         \STATE $type = null$
     \ENDIF
\ELSIF{$isA(exprNode, OperatorNode)$}
     \STATE $type = determineOperatorType(exprNode, contextVarTypes)$
\ELSIF{$isA(exprNode, FLWORExprNode)$}
     \STATE $returnClauseNode = ((FLWORExprNode)exprNode).returnClauseNode$
     \STATE $type = returnClauseNode.type$
\ENDIF

\RETURN $type$
\end{algorithmic}
\end{algorithm}

\todo[inline]{Napisat, ze determineOperatorType() tiez nie je kompletna}
\todo[inline]{definovat isOperatorClassAddition()}
\todo[inline]{definovat isNumericType()}
\todo[inline]{definovat isOperatorClassComparison()}

\begin{algorithm}
\caption{Function determineOperatorType}
\label{ALG_func_determineOperatorType}
\begin{algorithmic}[1]
\REQUIRE{\ \\
    $operatorNode$: An operator expression node in the syntax tree. \\
    $contextVarTypes$: Types of local variables valid in the current subtree.
}
\ENSURE The operator expression type.

\STATE $type = null$
\STATE $operator = operatorNode.operator$
\IF{$isOperatorClassComparison(operator)$}
    \STATE $type = xs:boolean$
\ELSIF{$isOperatorClassAddition(operator) AND$ one of the operands is of a numeric type}
    \STATE $leftOperand = ((OperatorNode)exprNode).leftOperand$
    \STATE $rightOperand = ((OperatorNode).exprNode).rightOperand$
    \IF{$isNumericType(leftOperand.type) AND isNumericType(rightOperand.type)$}
        \STATE $type = selectMoreGeneralType(leftOperand.type, rightOperand.type)$
    \ELSIF{$isNumericType(leftOperand.type)$}
        \STATE $type = leftOperand.type$
    \ELSIF{$isNumericType(rightOperand.type)$}
        \STATE $type = rightOperand.type$
    \ENDIF
\ENDIF

\RETURN $type$
\end{algorithmic}
\end{algorithm}

The algorithm iterates through the variable declaration nodes from the prolog and if the variable is explicitly assigned with its type, the type is noted. Otherwise, an attempt of the type deduction of the expression is made. The deduction of the type from the expression is presented in Algorithm \ref{ALG_func_determineExpressionType}. On input, it takes an expression node and types of local variables that are valid in the expression (also called context of variables or variable context). The presented algorithm is just shortened illustration since the complete version is too long to be presented in this text. Nevertheless, the principle is straightforward. Depending on the class of the expression, its type can be either determined directly or it depends on its subexpressions. \todo{link na mozne triedy vyrazov.}

Function $determineExpressionType()$ is called within function $analysisOfExpressionTypes$ presented in Algorithm \ref{ALG_func_analysisOfExpressionTypes}, which recursively determines types of all subexpressions. Two important questions about this function emerge. How (in which order) are the tree nodes recursively processed and how does the function handle definitions (bindings) of new variables which may also appear in the expressions?

The order of the recursion is post-order. The reason is that an expression node may need to know types of its subexpressions to determine its own type and we want to determine types of all expressions. Therefore, using the post-order, subexpressions of an expression are processed first, and then the expression itself without any need for further recursion, because the types of the subexpressions are already determined.

The handling of new variable definitions relies on the fact that the definitions make the new variables valid only in nodes with a higher sequence number when numbered in the post-order \todo{overit}. Therefore, the left-most subnode of a certain node can be processed without an extension of the variable context. And, if the left-most subnode extends the variable context for the following nodes in the post-order numbering, it can be easily handled because the types of the new variables can be directly determined since the binding expressions are already processed (thus, their types are known). Every node ''knows'' wheather or not it may define new variables for nodes with higher sequence numbers. This is expressed in the algorithm by the condition at line 10. This condition is double-checked at line 3 where its meaning is that every node even ''knows'' if its particular subnode may define new variables. An example is a $FLWOR$ construct. The processing of $FLWOR$ knows that the first part (the first subnode) of this node may bind new variables using $for$ and $let$ constructs. After the subnode is processed, the processing of the $FLWOR$ node looks if any new variables were bound and if so the context of local variables is extended. This is done at line 4. Function $mergeContextVarTypes$ presented in Algorithm \ref{ALG_function_mergeContextVarTypes} writes every record from its second argument (the new variables) into ist first argument (the variable context). If there are records for variables with the same names in the context, they are overwritten to correspond to the variable overlaping. It is important to note that we assume the variable context is not passed by refence but by value, and hence, the variable context of a certain node is not affected by the recursive processing of its subnodes.

\begin{algorithm}
\caption{Function analysisOfExpressionTypes}
\label{ALG_func_analysisOfExpressionTypes}
\begin{algorithmic}[1]
\REQUIRE{\ \\
    $startingNode$: A node determinig a subtree to perform the analysis on \\
    $contextVarTypes$: Types of local variables valid in the current subtree.
}

\FORALL{$subnode \in$ all subnodes of $startingNode$}
    \STATE $analysisOfExpressionTypes(subnode, contextVarTypes)$ \COMMENT{recursion}
    \IF{$subnode$ defines new variables valid for following subnodes}
        \STATE $contextVarTypes = mergeContextVarTypes(contextVarTypes, subnode.variableTypes)$
    \ENDIF
\ENDFOR
\IF{$startingNode$ is expression node}
    \STATE $startingNode.type = determineExpressionType(startingNode)$
\ENDIF
\IF{$startingNode$ defines new variables valid in following nodes}
    \STATE determine types of the variables
    \STATE memorize their types in $startingNode.variableTypes$
\ENDIF
\end{algorithmic}
\end{algorithm}

\todo[inline]{Algoritmus analysisOfEXpressionTypes, line 10 - Tady tyhle věty jsou v pseudoalgoritmech divné. Měy by se vyskytovat minimálně. JInak je to pak v podstatě jen odřádkovaná věta, ne formální popis algoritmu bez technických detailů.}

\todo[inline]{Pozor na to přetékání textu - na spoustě míst.}

\begin{algorithm}
\caption{Function mergeContextVarTypes}
\label{ALG_function_mergeContextVarTypes}
\begin{algorithmic}[1]
\REQUIRE{\ \\
    $contextVarTypes$: Types of local variables valid in current context. \\    
    $extendingVarTypes$: Variable types of the same structure as the first argument and extending it.
}
\ENSURE The context variable types extended with the variable types from the second argument.

\FORALL{$varName \in keys(extendingVarTypes)$}
    \STATE $contextVarTypes[varName] = extendingVarTypes[varName]$
\ENDFOR
\RETURN $contextVarTypes$
\end{algorithmic}
\end{algorithm}

Function $analysisOfExpressionTypes$ called upon a binding expression of a global variable determines the type of the binding expression, and hence, the type of the variable. In case of global variables, the argument $contextVarTypes$ is empty, because there are no local variables valid in the prolog section.

\begin{algorithm}
\caption{Function getVariableType}
\label{ALG_function_get_variable_type}
\begin{algorithmic}[1]
\REQUIRE{\ \\
    $contextVarTypes$: Types of local variables valid in current context. \\
    $var$: A variable to determine the type of.
}
\ENSURE The type of variable $var$.

\IF{$contextVarTypes[var]$ is defined}
     \RETURN $contextVarTypes[var]$
\ELSE
     \RETURN $getGlobalVariableType(var)$
\ENDIF
\end{algorithmic}
\end{algorithm}

The types of processed global variables are available trough function $getGlobalVariableType()$. If the function gets a variable that is not a global one, it returns $null$. Function $getVariableType()$ called in Algorithm \ref{ALG_func_determineExpressionType} is defined in Algorithm \ref{ALG_function_get_variable_type}. It checks if the specified variable is amongst the given local variables. If so, it returns its type, else it handles the variable as a global one and if such global variable does not exist the result is $null$.

\subsection{Determination of Expression Types}
To determine types of expressions, we use already introduced function $analysisOfExpressionTypes$. The starting node (its first argument) is the node representing the query body and the variable context is empty as there cannot be any local variable valid in the body node.

We can also determine expression types in functions. To do this for a certain function, $analysisOfExpressionTypes$ function has to be called with a function declaration node as the starting node. In this case, the function declaration node contains a subnode specifying function's formal arguments. These argument are set as the variable context for the function body represented by another subnode.

\section{Step 3: Inference of Built-in Types}
In this step, the algorithm goes through the syntax tree to infer types of elements and attributes from expressions using the type information from the previous step. These two steps could be merged together but for better comprehension we present it separately.

How are the types inferred from the expression types? We do not exploit all expressions. Only expressions of a particular type are exploited. Specifically, an expression has to contain an subexpression $E$ of $pathType$ type (expression representing a certain element or attribute or set of elements or attributes). In the following text, the set represented by expression $E$ is denoted $S$. Another requirement is that the expression has to be either a function call or an arithmetic operation \todo{Q: Nejake ine ako napr type constructor, cast as? Princip je stale rovnaky. A: ak to jen popište. To je naprosto OK.}.

The output of this step is a set of couples ($pathType$, XML Schema built-in type).

\todo[inline]{Zlucovanie odvodenych tvrdeni.}
\todo[inline]{Path expressions s predikatmi.}

\subsection{Function Calls}
This case is quite straightforward. The algorithm encounters a function call and one of the arguments is a known set of elements of attributes (subexpression $E$ representing $S$). To determine the type of $S$, it is only needed to determine the type of the corresponding formal argument from the definition of the function. The function is either a built-in one so its definition is known or it is defined in the prolog section. External functions are not processed as was mentioned already.

If the type $T$ of the formal argument is a built-in type or its sequence, then $T$ is also the inferred type of $S$. Otherwise, no statement is inferred.

\subsection{Arithmetic Operations}
If the operator is one of $+$, $-$, $*$, $div$, $mod$ (the class of the expression is $Operator$ and the class of the operator is one of $PLUS$, $MINUS$, $MUL$, $DIV$, $MOD$ \todo{Niekam napisat mozne triedy operatora}), one operand is $E$ and the type $T$ of the other operand is one of numeric built-in types, then $T$ is also the inferred type of $S$.

If the operator is one of $<$, $>$, $<=$, $>=$, $=$, $!=$ (the class of the expression is $Operator$ and the class of the operator is one of $GEN\_LESS\_THAN$, $GEN\_GREATER\_THAN$, $GEN\_LESS\_THAN\_EQUALS$, $GEN\_GREATER\_THAN\_EQUALS$, $GEN\_EQUALS$, $GEN\_NOT\_EQUALS$), one operand is $E$ and the type $T$ of the other operand is one of built-in types, then $T$ is also the inferred type of $S$.

\todo[inline]{Pěkné, neškodilo by proložit příklady, tj. dotaz, strom, co z něho odvodíme.}
\chapter{Combination with Existing Methods of Inference} \label{CHAPTER_combination_with_existings_methods_of_inference}
In this chapter, we describe how to incorporate the inferred statements to existing methods of XML schema inference. We focus on a class of methods which are based on a creation and subsequent simplification of the initial grammar (as discussed in Chapter \ref{chapter_analysis_of_recent_approaches}).

The initial grammar contains rules of form $e \rightarrow e_1e_2\dots e_k$, where $e$ is an element and $e_1, \dots , e_k$ are its subelements. After the simplification, the simplified grammar contain rules of form $e \rightarrow E$, where $E$ is a regular expression composed of subelements of element $e$, describing its content.

Attributes of elements are not contained in the grammar directly, but every element carries an information on its attributes.

The combination with such methods of inference is straightforward. The rules of the grammar describe a XML structure using the most general aspect of elements and their subelements. Since the statements inferred by our method do not involve the XML structure by defining a subelement structure of elements, there are no conflicts between the grammar rules and our statements that need to be resolved.

Our statements are of the three following forms. The first one is $p \rightarrow t$, where $p$ is a PathType representing an XQuery path selecting a set of elements or attributes, and $t$ is an XSD built-in atomic type.

The second and third forms are $k = (C,P,\{L\})$ and $fk = (C(P_1,\{L_1\}) \rightarrow (P_2,\{L_2\}))$, representing a key and a foreign key, where $C,P,P_1,P_2,L,L_1,L_2$ are PathTypes without predicates using only child and descendant axes.

In all tree cases, we want to determine elements from the grammar (or their attributes), targeted by the respective paths ($p$ from the first form and $C$ from the second and third form). This is done in two steps. The first one is normalization of a particular PathType and the second one is selection of the targeted elements.

\section{Evaluation of Paths}
A path represented by PathType can contain variable references and association of these references with other paths. That is the reason why the normalization is convenient. It simply finds the variable references in the path and replaces them with steps from the associated paths, which are normalized recursively.

The selection of elements is a simplified XPath evaluation. The simplification involves two aspects; the evaluation is not performed upon XML data, but the simplified grammar containing rules with a regular expression on their right side, and, partially related, predicates in path steps are not evaluated, they are ignored.

It iterate through steps of a path, maintaining a so-called context set, which is a set of elements to evaluate the current step upon. The evaluation of one step is shown in Algorithm \ref{FUNC_evaluate_step}. For the retaining of readability of the code, we present an evaluation of self or descendant, child, and attribute axes, and nodes specified by name.

\begin{algorithm}
\caption{Function evaluateStep}
\label{FUNC_evaluate_step}
\begin{algorithmic}[1]
\REQUIRE{\ \\
	$step$: An instance of StepExprNode representing a step of the path to evaluate.
	$contextSet$: A set of elements and attributes to evaluate $step$ upon.
	$grammar$: The grammar.
}

\ENSURE A context set after the step evaluation.

\IF{$is(step,$ SelfOrDescendantStep$)$}
	\RETURN $evaluateStep\_selfOrDescendant(contextSet, grammar)$
\ENDIF

\STATE $newContextSet :=$ an empty set

\STATE $axisKind := step.getChild(axisNode).axisKind$
\STATE $nodeName := step.getChild(axisNode).getChild(nameTestNode).name$
\IF{$axisKind =$ CHILD}
	\FORALL{$node \in contextSet$}
		\IF{$node$ is element}
			\STATE $subelements := getTokens(grammar[node])$
			\FORALL{$subelement \in subelements$}
				\IF{$subelement.name = nodeName$}
					\STATE add $subelement$ to $newContextSet$
				\ENDIF
			\ENDFOR
		\ENDIF
	\ENDFOR
\ELSIF{$axisKind =$ ATTRIBUTE}
	\FORALL{$node \in contextSet$}
		\IF{$node$ is element}
			\STATE $attributes := node.attributes$
			\FORALL{$attribute \in attributes$}
				\IF{$attribute.name = nodeName$}
					\STATE add $attribute$ to $newContextSet$
				\ENDIF
			\ENDFOR
		\ENDIF
	\ENDFOR
\ENDIF

\RETURN $newContextSet$
\end{algorithmic}
\end{algorithm}

At first, the algorithm determine the axis of the step. If it is self or descendant axis, it returns the result of \texttt{evaluateStep\_selfOrDescendant} function. This function returns the given context set extended by all descendants of the elements from the context set in the given grammar.

If the axis is child axis, for each element in the context set, the algorithm determines its subelement using the grammar, and, if those with the name equal to the name specified by to step are added to the resulting context set. Function \texttt{getTokens} at line 10 retrieves all elements from a regular expression. A specific form of the regular expression is not important, because we only need to know which elements are possible subelements of the particular element.

And, at last, if the axis is attribute axis, elements from the context set are searched to contain an attribute with the specified name and those found attributes are added to the result.

\section{Saving the Inferred Statements}
In the case of keys, the algorithm iterates through the inferred key statements. A key's context path $C$ is evaluated and the key is assigned to the target elements selected by the context path, one element can be assigned with multiple keys. Additionally, each key is assigned with a list of foreign keys that are referencing it.

In the case of inferred types, the situation is slightly less simple, because there can occur conflicting statements. The examples of such conflicts for path $p$ without predicates are $p_1 \rightarrow date$, $p_2 \rightarrow string$, and, $p_1 \rightarrow byte$, $p_2 \rightarrow int$, where $p_1$, $p_2$ are paths that when stripped of predicates, they equal $p$.

Note that, in both examples, one type is castable to the other (date to string, and byte to int). Both types are inferred correctly for nodes targeted by $p$, but one of them was inferred from a more convenient expression and is more precise.

Consider these two expressions; PathType $p$ is compared to an integer literal constant, and, PathType $p$ is an argument of function where a formal type of the argument is byte. The first expression is utilized to infer statement $p \rightarrow integer$, and the second one to infer statement $p \rightarrow byte$. Both of the types are correct, but the second one is more accurate.

A problem emerges for example if PathType $p$ is compared to an integer constant, and the real type of elements (or attributes) selected by $p$ is \texttt{double}. In that case, statement $p \rightarrow$ \texttt{integer} is inferred, but it is not correct.

\subsection{Verification using XML data}
To solves those problems, we propose a simple verification using XML data.

For each normalized PathType $p$ from the inferred type statements $\mathcal{S}_t$, we find set $T_p$ of all inferred types. $T_p = \{t|(p \rightarrow t) \in \mathcal{S}_t\}$. Then, we create sequence $T_p'$ by ordering the set $T_p$ from the most specific type to the most general one. For example, if $T_p = \{$\texttt{double}, \texttt{byte}, \texttt{int}$\}$, $T_p' = ($\texttt{byte}, \texttt{int}, \texttt{double}$)$.

Since we have the XML data and path $p$ is XQuery path, we can use an XQuery process (a program that evaluates XQuery paths or queries) to select nodes $N$ targeted by $p$. The verification algorithms iterates through $T_p'$ and for each $t' \in T_p'$ it checks if every node in $N$ conforms to $t'$. If so, $t'$ is the inferred type for nodes $N$ (and PathType $p$), else the inferred type is \texttt{string}.

\chapter{Implementation}
\chapter{Experiments}
\chapter{Conclusion}
The aim of this thesis was to employ XML operations in the XML schema inference. We analyzed several existing methods of the XML schema inference and we searched for methods that utilize some XML operations, only to found one\todo{Jiny vyraz}, which infers keys from XQuery queries.

Since the notion of XML operations is very general and the range of XML technologies is very large, for the purpose of this work, we decided to focus on XQuery technology. We made the overview of possible utilization of XQuery queries in the process of XML schema inference.

Before a creation of the algorithm itself, we had to take several decisions in questions that emerged. Since there is a lack of the methods dealing with the utilisation of XML operation, there is also a lack of practically proven solutions that can help in such decision making.

In the proposed solution, we decided to incorporate lexical and syntax analyses of XQuery queries, because it is more general and more extensible than a pattern searching. To achieve that, we adopted the algorithm from a recent master thesis dealing with an analysis of XQuery queries.

We also implemented several ideas from the overview to infer XSD built-in types of elements and attributes. And we extended the one existing method dealing with the inference of keys. We experimentally demonstrated that on some input files, results of the extended method are better than results of the original method. However, we did not succeed in a search of an ideal, large enough set of test data. The testing was performed on a small set of input queries and further testing and algorithm tuning is required.

Finally, we proposed a simple way how to combine the inferred statements with existing grammar inferring methods of the XML schema inference and we implemented it using the jInfer framework. Thus, we created the first complete, ready-to-use, and extensible implementation of the XML xchema inference exploiting XQuery queries besides XML data.

\todo[inline]{Nejake prehladne zhrnutie plusov a minusov?}

\section{Future Work}
This work implements only some ideas discussed in the overview in Chapter \ref{CHAPTER_analysis_of_xquery}, leaving most of them for a future work. Also, the implemented algorithms can be further refined as was already mentioned in Chapter \ref{CHAPTER_proposed_algorithm} with the presentation of the algorithms. For example thorough processing of user-defined function calls.

Chapter \ref{CHAPTER_precreation_of_algorithm} discusses some possible future enhancements as well. A research on possibilities of modifying the grammar rules based on information extracted from XQuery queries may bring interesting results. We analyse the queries statically only, we do not evaluate them. An analysis of queries together with their results can be a topic of another future research direction.

The utilization of XQuery queries certainly provides a space for incorporating an interaction with user. For example, a user may influence the scoring of inferred keys to get more precise results.

And, a very large space for a possible future research is provided by utilisation of other XML operations. The main representant is XSLT.

Besides the mentioned, our opinion is that the most urgent future work is obtaining a large enough test data set, performing proper experiments, and refining the algorithm by modification of its settings (statements inferred from join pattern occurrences, weights, etc.) according to experimental results. This process was suggested in Chapter \ref{CHAPTER_experiments}.

%%% Seznam pouľité literatury
\nocite{*}
\bibliographystyle{plain}
\bibliography{literature}

%%% Pouľité zkratky v diplomové práci, existují-li, včetně jejich vysvětlení.
%\chapwithtoc{List of Abbreviations}
%\todo[inline]{Není nutno.}

%%% Přílohy k diplomové práci, existují-li (různé dodatky jako výpisy programů,
%%% diagramy apod.). Kaľdá příloha musí být alespoň jednou odkazována z vlastního
%%% textu práce. Přílohy se číslují.
%\chapwithtoc{Attachments}
\appendix
\chapter{Attachments}
\begin{figure}
\caption{All possible values representing an operator in an instance of \texttt{OperatorNode}}
\label{FIG_operators}
Logical: \texttt{AND}, \texttt{OR}.

Comparison: \texttt{GEN\_EQUALS}, \texttt{GEN\_NOT\_EQUALS}, \texttt{GEN\_LESS\_THAN}, \texttt{GEN\_LESS\_THAN\_EQUALS}, \texttt{GEN\_GREATER\_THAN}, \texttt{GEN\_GREATER\_THAN\_EQUALS}, \texttt{VAL\_EQUALS}, \texttt{VAL\_NOT\_EQUALS}, \texttt{VAL\_LESS\_THAN}, \texttt{VAL\_LESS\_THAN\_EQUALS}, \texttt{VAL\_GREATER\_THAN}, \texttt{VAL\_GREATER\_THAN\_EQUALS}, \texttt{NOD\_IS}, \texttt{NOD\_PRECEDES}, \texttt{NOD\_FOLLOWS}.

Range: \texttt{TO}.

Additive: \texttt{PLUS}, \texttt{MINUS}.

Multiplicative: \texttt{MUL}, \texttt{DIV}, \texttt{IDIV}, \texttt{MOD}.

Set: \texttt{UNION}, \texttt{INTERSECTION}, \texttt{DIFFERENCE}.

Type test: \texttt{INSTANCE\_OF}, \texttt{CASTABLE\_AS}.

Type conversion: \texttt{TREAT\_AS}, \texttt{CAST\_AS}.

Unary: \texttt{UNARY\_PLUS}, \texttt{UNARY\_MINUS}.
\end{figure}

\chapter{Test Data} \label{APPENDIX_test_data}

Test data from the XMark project \cite{xmark}. Using the provided XML generator, XML document of size approximately 1.5 MB was generated. Figure \ref{listing_dtd} is its DTD and other figures in this appendix are XQuery queries that query the XML document.

The queries was slightly modified by replacing calls of \texttt{doc} function in paths by \texttt{/} (document node).

\begin{lstlisting}[caption=DTD of the test XML data, frame=single, label=listing_dtd]
<!-- DTD for auction database -->
<!-- $Id: auction.dtd,v 1.15 2001/01/29 21:42:35 albrecht Exp $ -->

<!ELEMENT site            (regions, categories, catgraph, people, open_auctions, closed_auctions)>

<!ELEMENT categories      (category+)>
<!ELEMENT category        (name, description)>
<!ATTLIST category        id ID #REQUIRED>
<!ELEMENT name            (#PCDATA)>
<!ELEMENT description     (text | parlist)>
<!ELEMENT text            (#PCDATA | bold | keyword | emph)*>
<!ELEMENT bold            (#PCDATA | bold | keyword | emph)*>
<!ELEMENT keyword         (#PCDATA | bold | keyword | emph)*>
<!ELEMENT emph            (#PCDATA | bold | keyword | emph)*>
<!ELEMENT parlist         (listitem)*>
<!ELEMENT listitem        (text | parlist)*>

<!ELEMENT catgraph        (edge*)>
<!ELEMENT edge            EMPTY>
<!ATTLIST edge            from IDREF #REQUIRED to IDREF #REQUIRED>

<!ELEMENT regions         (africa, asia, australia, europe, namerica, samerica)>
<!ELEMENT africa          (item*)>
<!ELEMENT asia            (item*)>
<!ELEMENT australia       (item*)>
<!ELEMENT namerica        (item*)>
<!ELEMENT samerica        (item*)>
<!ELEMENT europe          (item*)>
<!ELEMENT item            (location, quantity, name, payment, description, shipping, incategory+, mailbox)>
<!ATTLIST item            id ID #REQUIRED
                          featured CDATA #IMPLIED>
<!ELEMENT location        (#PCDATA)>
<!ELEMENT quantity        (#PCDATA)>
<!ELEMENT payment         (#PCDATA)>
<!ELEMENT shipping        (#PCDATA)>
<!ELEMENT reserve         (#PCDATA)>
<!ELEMENT incategory      EMPTY>
<!ATTLIST incategory      category IDREF #REQUIRED>
<!ELEMENT mailbox         (mail*)>
<!ELEMENT mail            (from, to, date, text)>
<!ELEMENT from            (#PCDATA)>
<!ELEMENT to              (#PCDATA)>
<!ELEMENT date            (#PCDATA)>
<!ELEMENT itemref         EMPTY>
<!ATTLIST itemref         item IDREF #REQUIRED>
<!ELEMENT personref       EMPTY>
<!ATTLIST personref       person IDREF #REQUIRED>

<!ELEMENT people          (person*)>
<!ELEMENT person          (name, emailaddress, phone?, address?, homepage?, creditcard?, profile?, watches?)>
<!ATTLIST person          id ID #REQUIRED>
<!ELEMENT emailaddress    (#PCDATA)>
<!ELEMENT phone           (#PCDATA)>
<!ELEMENT address         (street, city, country, province?, zipcode)>
<!ELEMENT street          (#PCDATA)>
<!ELEMENT city            (#PCDATA)>
<!ELEMENT province        (#PCDATA)>
<!ELEMENT zipcode         (#PCDATA)>
<!ELEMENT country         (#PCDATA)>
<!ELEMENT homepage        (#PCDATA)>
<!ELEMENT creditcard      (#PCDATA)>
<!ELEMENT profile         (interest*, education?, gender?, business, age?)>
<!ATTLIST profile         income CDATA #IMPLIED>
<!ELEMENT interest        EMPTY>
<!ATTLIST interest        category IDREF #REQUIRED>
<!ELEMENT education       (#PCDATA)>
<!ELEMENT income          (#PCDATA)>
<!ELEMENT gender          (#PCDATA)>
<!ELEMENT business        (#PCDATA)>
<!ELEMENT age             (#PCDATA)>
<!ELEMENT watches         (watch*)>
<!ELEMENT watch           EMPTY>
<!ATTLIST watch           open_auction IDREF #REQUIRED>

<!ELEMENT open_auctions   (open_auction*)>
<!ELEMENT open_auction    (initial, reserve?, bidder*, current, privacy?, itemref, seller, annotation, quantity, type, interval)>
<!ATTLIST open_auction    id ID #REQUIRED>
<!ELEMENT privacy         (#PCDATA)>
<!ELEMENT initial         (#PCDATA)>
<!ELEMENT bidder          (date, time, personref, increase)>
<!ELEMENT seller          EMPTY>
<!ATTLIST seller          person IDREF #REQUIRED>
<!ELEMENT current         (#PCDATA)>
<!ELEMENT increase        (#PCDATA)>
<!ELEMENT type            (#PCDATA)>
<!ELEMENT interval        (start, end)>
<!ELEMENT start           (#PCDATA)>
<!ELEMENT end             (#PCDATA)>
<!ELEMENT time            (#PCDATA)>
<!ELEMENT status          (#PCDATA)>
<!ELEMENT amount          (#PCDATA)>

<!ELEMENT closed_auctions (closed_auction*)>
<!ELEMENT closed_auction  (seller, buyer, itemref, price, date, quantity, type, annotation?)>
<!ELEMENT buyer           EMPTY>
<!ATTLIST buyer           person IDREF #REQUIRED>
<!ELEMENT price           (#PCDATA)>
<!ELEMENT annotation      (author, description?, happiness)>

<!ELEMENT author          EMPTY>
<!ATTLIST author          person IDREF #REQUIRED>
<!ELEMENT happiness       (#PCDATA)>
\end{lstlisting}

\begin{lstlisting}[caption=Test query 1., frame=single, label=listing_test_query_1]
for $b in /site/people/person[@id = "person0"] return $b/name/text()
\end{lstlisting}

\begin{lstlisting}[caption=Test query 2., frame=single, label=listing_test_query_2]
for $b in /site/open_auctions/open_auction
return <increase>{$b/bidder[1]/increase/text()}</increase>
\end{lstlisting}

\begin{lstlisting}[caption=Test query 3., frame=single, label=listing_test_query_3]
for $b in /site/open_auctions/open_auction
where zero-or-one($b/bidder[1]/increase/text()) * 2 <= $b/bidder[last()]/increase/text()
return
  <increase
  first="{$b/bidder[1]/increase/text()}"
  last="{$b/bidder[last()]/increase/text()}"/>
\end{lstlisting}

\begin{lstlisting}[float=htpb, caption=Test query 4., frame=single, label=listing_test_query_4]
for $b in /site/open_auctions/open_auction
where
  some $pr1 in $b/bidder/personref[@person = "person20"],
       $pr2 in $b/bidder/personref[@person = "person51"]
  satisfies $pr1 << $pr2
return <history>{$b/reserve/text()}</history>
\end{lstlisting}

\begin{lstlisting}[float=htpb, caption=Test query 5., frame=single, label=listing_test_query_5]
count(
  for $i in /site/closed_auctions/closed_auction
  where $i/price/text() >= 40
  return $i/price
)
\end{lstlisting}

\begin{lstlisting}[float=htpb, caption=Test query 6., frame=single, label=listing_test_query_6]
for $b in //site/regions return count($b//item)
\end{lstlisting}

\begin{lstlisting}[float=htpb, caption=Test query 7., frame=single, label=listing_test_query_7]
for $p in /site
return
  count($p//description) + count($p//annotation) + count($p//emailaddress)
\end{lstlisting}

\begin{lstlisting}[float=htpb, caption=Test query 8., frame=single, label=listing_test_query_8]
for $p in /site/people/person
let $a :=
  for $t in /site/closed_auctions/closed_auction
  where $t/buyer/@person = $p/@id
  return $t
return <item person="{$p/name/text()}">{count($a)}</item>
\end{lstlisting}

\begin{lstlisting}[float=htpb, caption=Test query 9., frame=single, label=listing_test_query_9]
let $ca := /site/closed_auctions/closed_auction return
let
    $ei := /site/regions/europe/item
for $p in /site/people/person
let $a :=
  for $t in $ca
  where $p/@id = $t/buyer/@person
  return
    let $n := for $t2 in $ei where $t/itemref/@item = $t2/@id return $t2
    return <item>{$n/name/text()}</item>
return <person name="{$p/name/text()}">{$a}</person>
\end{lstlisting}

\begin{lstlisting}[float=htpb, caption=Test query 10., frame=single, label=listing_test_query_10]
for $i in
  distinct-values(/site/people/person/profile/interest/@category)
let $p :=
  for $t in /site/people/person
  where $t/profile/interest/@category = $i
  return
    <personne>
      <statistiques>
        <sexe>{$t/profile/gender/text()}</sexe>
        <age>{$t/profile/age/text()}</age>
        <education>{$t/profile/education/text()}</education>
        <revenu>{fn:data($t/profile/@income)}</revenu>
      </statistiques>
      <coordonnees>
        <nom>{$t/name/text()}</nom>
        <rue>{$t/address/street/text()}</rue>
        <ville>{$t/address/city/text()}</ville>
        <pays>{$t/address/country/text()}</pays>
        <reseau>
          <courrier>{$t/emailaddress/text()}</courrier>
          <pagePerso>{$t/homepage/text()}</pagePerso>
        </reseau>
      </coordonnees>
      <cartePaiement>{$t/creditcard/text()}</cartePaiement>
    </personne>
return <categorie>{<id>{$i}</id>, $p}</categorie>
\end{lstlisting}

\begin{lstlisting}[float=htpb, caption=Test query 11., frame=single, label=listing_test_query_11]
for $p in /site/people/person
let $l :=
  for $i in /site/open_auctions/open_auction/initial
  where $p/profile/@income > 5000 * exactly-one($i/text())
  return $i
return <items name="{$p/name/text()}">{count($l)}</items>
\end{lstlisting}

\begin{lstlisting}[float=htpb, caption=Test query 12., frame=single, label=listing_test_query_12]
for $p in /site/people/person
let $l :=
  for $i in /site/open_auctions/open_auction/initial
  where $p/profile/@income > 5000 * exactly-one($i/text())
  return $i
where $p/profile/@income > 50000
return <items person="{$p/profile/@income}">{count($l)}</items>
\end{lstlisting}

\begin{lstlisting}[float=htpb, caption=Test query 13., frame=single, label=listing_test_query_13]
for $i in /site/regions/australia/item
return <item name="{$i/name/text()}">{$i/description}</item>
\end{lstlisting}

\begin{lstlisting}[float=htpb, caption=Testing query 14., frame=single, label=listing_test_query_14]
for $i in /site//item
where contains(string(exactly-one($i/description)), "gold")
return $i/name/text()
\end{lstlisting}

\begin{lstlisting}[float=htpb, caption=Test query 15., frame=single, label=listing_test_query_15]
for $a in
  /site/closed_auctions/closed_auction/annotation/description/parlist/
   listitem/
   parlist/
   listitem/
   text/
   emph/
   keyword/
   text()
return <text>{$a}</text>
\end{lstlisting}

\begin{lstlisting}[float=htpb, caption=Test query 16., frame=single, label=listing_test_query_16]
for $a in /site/closed_auctions/closed_auction
where
  not(
    empty(
      $a/annotation/description/parlist/listitem/parlist/listitem/text/emph/
       keyword/
       text()
    )
  )
return <person id="{$a/seller/@person}"/>
\end{lstlisting}

\begin{lstlisting}[float=htpb, caption=Test query 17., frame=single, label=listing_test_query_17]
for $p in /site/people/person
where empty($p/homepage/text())
return <person name="{$p/name/text()}"/>
\end{lstlisting}

\begin{lstlisting}[float=htpb, caption=Test query 18., frame=single, label=listing_test_query_18]
declare namespace local = "http://www.foobar.org";
declare function local:convert($v as xs:decimal?) as xs:decimal?
{
  2.20371 * $v (: convert Dfl to Euro :)
};

for $i in /site/open_auctions/open_auction
return local:convert(zero-or-one($i/reserve))
\end{lstlisting}

\begin{lstlisting}[float=htpb, caption=Test query 19., frame=single, label=listing_test_query_19]
for $b in /site/regions//item
let $k := $b/name/text()
order by zero-or-one($b/location) ascending empty greatest
return <item name="{$k}">{$b/location/text()}</item>
\end{lstlisting}

\begin{lstlisting}[float=htpb, caption=Test query 20., frame=single, label=listing_test_query_20]
<result>
  <preferred>
    {count(/site/people/person/profile[@income >= 100000])}
  </preferred>
  <standard>
    {
      count(
        /site/people/person/
         profile[@income < 100000 and @income >= 30000]
      )
    }
  </standard>
  <challenge>
    {count(/site/people/person/profile[@income < 30000])}
  </challenge>
  <na>
    {
      count(
        for $p in /site/people/person
        where empty($p/profile/@income)
        return $p
      )
    }
  </na>
</result>
\end{lstlisting}
\chapter{Content of CD}
The CD attached to this thesis has the following structure.

\begin{itemize}
\item \texttt{content.txt} - A file with this text.
\item \texttt{text/} - A PDF version of the thesis.
\item \texttt{src/} - Source codes of the jInfer framework including implementation of our method. The same source codes can be also obtained from public Subversion repository by issuing command:\\
\texttt{svn co -r 2155 \\ https://jinfer.svn.sourceforge.net/svnroot/jinfer/jinfer/trunk/}
\item \texttt{bin/} - jInfer plugins for NetBeans 7.0.1. See jInfer Tutorial \cite{jinfer_tutorial} for step-by-step instructions, but use NetBeans 7.0.1. Since the tutorial is for the last official jInfer release 1.0, it says the required version of NetBeans is at least 6.9. But this is not true in our case, because we use the development version of jInfer and it requires NetBeans 7.0.1.
\item \texttt{testing/} - A directory containing sets of test data and test results. Again, see the Jinfer Tutorial \cite{jinfer_tutorial} for instructions how to run the inference. 
\end{itemize}

%%% Tabulky v diplomové práci, existují-li.
%\chapwithtoc{List of Tables}
\listoftables
\listoffigures

\openright
\end{document}
